%% A living glossary for the combinator-S series, S1--S12 and the ancillary SIM paper, with
%% the working labels of the unpublished Lock 1 notes in a separate part.  Started 2026-11-26.
%% Build: `make` (latexmk); read main.log for undefined references and overfull boxes.
\documentclass[11pt]{amsart}
\usepackage[T1]{fontenc}
\usepackage{lmodern}
\usepackage{amsmath,amssymb,amsthm}
\usepackage[numbers,sort]{natbib}
\usepackage{longtable}
\usepackage{array}
\usepackage[hidelinks]{hyperref}
\usepackage{bm}

\raggedbottom

\newcommand{\Ter}{\mathcal S}
\newcommand{\red}{\twoheadrightarrow}
\newcommand{\hred}{\twoheadrightarrow_h}
\newcommand{\conv}{=_\beta}
\newcommand{\sz}[1]{\lvert#1\rvert}
\newcommand{\rest}{\mathrm{rest}}
\newcommand{\nf}{\mathrm{nf}}
\newcommand{\fold}{\mathrm{fold}}
\newcommand{\pop}{\mathrm{pop}}
\newcommand{\Hc}{\mathcal H}
\newcommand{\Ic}{\mathcal I}
\newcommand{\Gc}{\mathcal G}
\newcommand{\Kc}{\mathcal K}
\newcommand{\Cc}{\mathcal C}
\newcommand{\Zc}{\mathcal Z}
\newcommand{\Self}{\mathrm{SELF}}
\newcommand{\Bk}{\mathfrak B}
\newcommand{\sqs}{\sqsubseteq}
\newcommand{\psq}{\sqsubset}

%% ---------------------------------------------------------------------------------------------
%% HOW TO ADD AN ENTRY
%%
%%   \begin{entry}{g:slug}{Headword}{YYYY-MM-DD}
%%   \src{S12, \S1}
%%   \lean{\texttt{Module.name}}
%%   One to three sentences, paraphrased from the source.  \see{g:other}
%%   \end{entry}
%%
%% Placement.  Entries are alphabetical by headword within each part.  Headwords that start with a
%% numeral come first.  A symbol sorts by the name of its letter ($\beta$ under "beta", $\Hc$ under "H").
%% Slugs: g:<word> for terms, g:sym-<letter> for symbol entries, g:wl-<label> for Part II.
%%
%% Slots, in this order, each on its own line: \src{...}; in Part II \status{...} and, if a paper
%% uses the label, \appears{...}; then \lean{...} (optional, only for a declaration that exists in
%% lean-ski); then the text.  \see{g:slug} (optional) cross-references another entry.
%% Location labels: "S6, Def.~3" or "S11, \S2"; copy the number from the paper's main.aux.  For notes:
%% "LOCK1-CT \S24".  A new paper: add its bib entry to refs.bib and a line to the list in the
%% subsection "The papers".
%%
%% Several meanings: replace \src by one \sense{S6, Def.~3} per meaning; put \lean at the end of its
%% sense.  For an overloaded symbol, write "Overloaded." as the first line of the entry.
%%
%% Status grammar (Part II only): \status{<state> [tag]; <conditions>}, where <state> is one of
%% proved, conditional, open, refuted, computed, and <tag> is [P], [C], [K], [Conj], [Sketch] or [Lit] (see the entry
%% "status tags").  Conditions name their inputs, e.g. "given (S1), T10 [K], T20 [C]".  A pure
%% definition has no \status.  Repair codes such as W26.3 name repair lines of a review in the notes.
%% "refuted [P]" means shown false by a proof; "refuted [C]" means a computed counterexample.
%%
%% Style.  A term entry opens with its genus ("A <kind> that ..."); a result entry opens with the
%% result word ("Lemma.", "Theorem.", "Conjecture.").  Present tense, active voice, one to three
%% sentences (per sense in a multi-sense entry).  Write a name or a count instead of "some",
%% "several", "certain" or "the data".  Make no claim about the current state of another file.
%%
%% Dates.  The date is the day the entry was added or its content last changed; a typo fix does not
%% move it.
%%
%% Notation.  A symbol also gets a row \sym{symbol}{meaning}{source}{g:slug} in the notation table:
%% the symbol in math mode, a meaning of at most eight words, the source label, and the slug of the
%% explaining entry.  Rows for Part II labels go below the "Working labels" line.
%% ---------------------------------------------------------------------------------------------
\makeatletter
\newcounter{sense}
\newenvironment{entry}[3]{%
  \par\addvspace{\medskipamount}\noindent\phantomsection
  \def\@currentlabel{#2}\label{#1}\setcounter{sense}{0}\raggedright\emergencystretch=3em%
  \textbf{#2}\hfill{\footnotesize\textsf{#3}}\par\nopagebreak\noindent\ignorespaces}%
  {\par}
\makeatother
\newcommand{\src}[1]{\textsf{Source:} #1.\enskip}
\newcommand{\lean}[1]{\textsf{Lean:} #1.\enskip}
\newcommand{\status}[1]{\textsf{Status:} #1.\enskip}
\newcommand{\appears}[1]{\textsf{In papers:} #1.\enskip}
\newcommand{\sense}[1]{\par\noindent\stepcounter{sense}(\arabic{sense})~[#1]\enskip}
\renewcommand{\see}[1]{$\rightarrow$\,\emph{\ref{#1}}}
\newcommand{\sym}[4]{#1 & #2 & #3 & \pageref{#4}\\}

\title[A Glossary for the Combinator $S$ Series]{A Glossary for the Combinator $\bm{S}$ Series}
\author{Claude Opus \& Fable with N.\ E.\ Davis}
\address{ZeroAttest, Champaign, Illinois, USA}
\email{NDavis@ZeroAttest.com}
\urladdr{https://orcid.org/0000-0002-1131-8319}
\date{\today}

\begin{document}
\begin{abstract}
This glossary lists the terms, symbols, named results and conjectures of the papers on the
combinator $S$. Each entry names the paper and the definition or section that introduces the item,
gives a short definition, and cites the Lean declaration, if one exists. Part~II records the working
labels of the unpublished Lock~1 notes and the current status of each. The glossary grows with the
series.
\end{abstract}
\maketitle

\section*{About this glossary}

\subsection*{Purpose}
This glossary fixes the vocabulary of the series. The papers reuse one another's terms, and some
words and letters mean different things in different papers. Each entry names the paper that
defines a term, gives the definition or section number, and paraphrases the definition in one or two
sentences. An entry does not replace the paper. When a word has several meanings, its entry lists
them as numbered senses; when a symbol is overloaded, the entry says so.

\subsection*{How to cite an entry}
For the mathematics, cite the paper that defines the term, not the glossary. To cite the glossary
itself, give the headword and the date at the right of the entry, for example ``Glossary, entry
\emph{stable frame} (2026-11-26)''.

\subsection*{Dating}
Entries are added over time. The date at the right of an entry is the day it was added or last
changed; any change of content moves the date. A status in Part~II holds on that date only and may
have changed since.

\subsection*{Labels}
``S6, Def.~3'' names a numbered environment of that paper; all theorem-like environments of a paper
share one counter. ``S11, \S2'' names a section. ``Lean'' names a declaration of
\texttt{lean-ski}~\citep{leanski}, written in Lean~4~\citep{lean4}.

\subsection*{The papers}
S1~\citep{s1} head normal forms; S2~\citep{s2} productive terms; S3~\citep{s3} infinitary normal
forms; S4~\citep{s4} join bounds and the word problem; S5~\citep{s5} depth-two prefixes;
S6~\citep{s6} prefix regularity; S7~\citep{s7} no ground loops; S8~\citep{s8} plain contexts;
S9~\citep{s9} consumption; S10~\citep{s10} clock closure; S11~\citep{s11} one term determined by its teeth; S12~\citep{s12} counter
runs; S14~\citep{s14} a family determined by its trees; and the ancillary SIM~\citep{simech}, which mechanizes the self-interpreter of~\citep{davis-si}.
Where its usage differs, an entry cites an earlier draft of S8 on combs as
``s8-combs (draft)''. The background sources are \emph{Dance of the Starlings}~\citep{dots},
Waldmann's work~\citep{waldmann-thesis,waldmann-ic}, and the standard
references~\citep{barendregt,terese}.

{\raggedright Part~II cites the notes by file name: LOCK1-CT, LOCK1-CTB, LOCK1-structure,
LOCK1-NTD-probe, LOCK1-INJ, LOCK1-INV, LOCK1-CONG, LOCK1-FF, LOCK1-FFALL, A1-lock1-S3 and OPEN-QUESTIONS. It marks status with the tags explained
in the entry \emph{status tags}. The notes are unpublished; a label from them is a working name, not
a published result.\par}

\section*{Notation}

Each row gives a symbol, a short meaning, the defining paper or note, and the page of the entry that
explains it. Overloaded letters have several rows or a list of papers.

\begingroup\small
\begin{longtable}{@{}>{\raggedright}p{0.2\textwidth}>{\raggedright}p{0.4\textwidth}%
>{\raggedright}p{0.24\textwidth}l@{}}
Symbol & Meaning & Source & Page\\ \hline\endhead
\sym{$\Ter$}{closed terms built from $S$}{S1, \S2}{g:closed-s-terms}
\sym{$\Ter^{\infty}$}{finite and infinite $S$-trees}{S3, \S2}{g:closed-s-terms}
\sym{$\to$, $\red$}{weak reduction, its closure}{S1, \S2}{g:weak-reduction}
\sym{$\to_h$, $\hred$}{head step, head reduction}{S2, \S2; S11, \S1}{g:head-step}
\sym{$\conv$}{convertibility}{S4, \S1}{g:conv}
\sym{$\sz{t}$}{number of occurrences of $S$}{S2, \S2}{g:size}
\sym{$\nf(t)$}{normal form of $t$}{S12, \S1}{g:normalizes}
\sym{$h(t)$}{hnf reached by head reduction}{S2, \S2}{g:sym-h}
\sym{$\operatorname{fr}(t)$}{spine (list of items) of $t$}{S1, \S2}{g:spine}
\sym{$[u_1,\dots,u_m]$}{the frame $S\,u_1\cdots u_m$}{S8, \S1}{g:frame}
\sym{$S^m\,t$}{$m$ nested copies of $S$}{S12, \S1}{g:sym-Sm}
\sym{$\tau_n(t)$}{$n$-th tooth}{S11, \S1}{g:tooth}
\sym{$\rest(t)$}{function part of first stable frame}{S11, \S1}{g:rest}
\sym{$\sqsubseteq$, $\sqsubset$}{right suffix, proper right suffix}{S12, \S2}{g:right-suffix}
\sym{$t^{\infty}$}{infinitary normal form}{S3, Thm.~7}{g:t-infty}
\sym{$\mathrm{BT}(t)$}{B\"ohm tree}{S2, \S2}{g:bohm-tree}
\sym{$\mathrm{BeT}(t)$}{Berarducci tree}{S3, \S4}{g:berarducci-tree}
\sym{$\mathrm{Fin}(t)$}{boundedly many root steps}{S3, \S3}{g:fin}
\sym{$\mathrm{NF}(d,t,P)$}{depth-$d$ prefix of $t^{\infty}$}{S3, Rem.~8}{g:depth-prefix}
\sym{$\pi'_d(t)$, $\pi_d(t)$}{depth-$d$ prefix, head $S$ shown}{S5, Def.~1; S6, \S2}{g:depth-prefix}
\sym{$\bullet$}{hole in a prefix}{S5, Def.~1}{g:hole}
\sym{$\bar t$}{fold into an $S$-model}{S5, \S2; S6, \S2}{g:fold}
\sym{$b(t)$, $\fold(t)$}{block of $t$}{S1, \S2; S12, \S5}{g:fold}
\sym{$c(t)$}{class of $t$}{S2, \S3}{g:fold}
\sym{$\operatorname{st}(t)$}{product state (class, block)}{S2, \S3}{g:product-state}
\sym{$O(t)$}{observable over a model}{S6, Def.~3}{g:observable}
\sym{$\Psi(t)$}{values of first stable frame}{S6, Def.~3}{g:sym-Psi}
\sym{$M_d$, $o_d$}{finite model for depth-$d$ prefixes}{S6, Thm.~1}{g:prefix-regularity}
\sym{$\simeq_g$}{finite-index term equivalence}{S6, Thm.~9}{g:finite-index}
\sym{$\mathrm{rd}$}{right depth}{S7, \S2}{g:right-depth}
\sym{$a\cdot_k b$}{labelled application}{S7, \S2}{g:labelled-reduction}
\sym{$\red_f$}{reduction by f-steps}{S7, \S4}{g:f-step}
\sym{$C[u]$, $E[p]$}{context filled with a term}{S7, \S4; S8, \S1}{g:context}
\sym{$M\Downarrow_n N$, $M\Uparrow$}{whnf in $n$ steps; none}{SIM, \S1}{g:whnf}
\sym{$\langle t\rangle$}{Scott encoding}{SIM, \S1}{g:scott-encoding}
\sym{$A$}{$S\,S\,S$; hypothesis (A); a set}{S2, \S1; S4, \S4}{g:sym-A}
\sym{$\mathbf B$, $B$, $\mathbb M$}{the 38 blocks}{S1, \S2; S6; S9; S12}{g:block-model}
\sym{$c$, $c_j$}{weights; class; constant; counter pieces}{S1; S2; S4; S11}{g:sym-c}
\sym{$D$, $d$, $D(v)$}{373-model; state; metric; $S\,v\,(S\,v)$}{S5; S1; S3; S12}{g:sym-D}
\sym{$E$}{exempt blocks $\{0,17\}$; a context}{S1, \S4; S8, \S1}{g:sym-E}
\sym{$G$, $G_e$, $G_1$, $G_2$}{$S\,Q\,Q$; heaps; graphs}{S11; S12; S6}{g:sym-G}
\sym{$\Gc$}{counter grammar $\Zc\times\Hc$}{S12, Def.~1}{g:counter-grammar}
\sym{$\Hc$, $\Ic$, $\Zc$}{parts of the counter grammar}{S12, Def.~1}{g:counter-grammar}
\sym{$\Kc$}{normal $c$ with $c\,W$ normalizing}{S12, \S1}{g:set-K}
\sym{$P$, $P_k$, $\mathcal P$}{$S\,S\,S$; approximants; right combs}{S2; S3; S4; S6}{g:sym-P}
\sym{$p_K$}{$S\,(S\,S)\,G\,(Q\,(Q\,G))\,S$}{S11, Thm.~1}{g:pK}
\sym{$a_m$, $b_m$, $i_m$}{spine and counter of the run of $p_K$}{S11, \S3}{g:pk-run}
\sym{$Q$, $\mathcal Q$, $Q_k$}{overloaded, see entry}{S1; S3; S4; S6; S9; S11}{g:sym-Q}
\sym{$r$, $R$}{weight; rejecting class; $S\,S$}{S1; S2; S4}{g:sym-r}
\sym{$T$, $T_E$}{$S\,S$ (S2, S3); $S\,S\,S$ (S12); a tree}{S2; S8; S12}{g:sym-T}
\sym{$U$, $U_j$}{$S\,U$ (S3); passive argument (S11); $S\,T\,T$ (S12)}{S3; S11; S12}{g:sym-U}
\sym{$X$}{overloaded; $S\,T\,S$ in S12}{S1--S3; S8; S12}{g:sym-X}
\sym{$Y$, $Y_n$}{teeth of $p_K$; $X$ or $S\,X\,c$}{S8; S11; S12}{g:sym-Y}
\sym{$\sigma_j$, $\Cc$}{combs $S$, $S\,S$, $S\,(S\,S)$, \dots}{S12, \S1}{g:comb}
\sym{$(A,W)$}{normal pair: stack, heap}{S12, \S2}{g:normal-pair}
\sym{$N$}{normal form of a tooth}{S12, Lemma~6}{g:new-state}
\sym{$\mathrm{pop}(Z)$}{two-item core of the first item}{S12, Lemma~15}{g:pop}
\sym{$\Bk$}{the set of S11, \S5}{S11, \S5}{g:frakB}
\sym{$\beta$}{block table; subscript of $\conv$}{S1, \S2; S4}{g:sym-beta}
\sym{$\bot$}{block 36; tree bottom}{S1, \S2; S2; S3}{g:sym-bot}
\sym{$\gamma$, $\mathsf r$}{39-class table; rejecting class 37}{S2, \S3}{g:sym-gamma}
\sym{$\delta$, $\circledast$}{35-state automaton; rejecting state}{S1, \S3}{g:automaton-35}
\sym{$\kappa(x)$}{depth along first items to no hnf}{S11, \S3}{g:kappa}
\sym{$\varphi$}{blocks to states; model to blocks}{S1, eq.~(2); S5}{g:sym-phi}
\sym{$\pi$}{potential}{S1, \S4}{g:sym-pi}
\sym{$\tau(t)$, $\#_S(t)$}{weight tie-breakers}{S1, Thm.~9}{g:sym-tau}
\sym{$\varpi$}{bound on non-free steps}{S6, Lemma~6}{g:productive-runs}
\sym{$\Self$}{self-dead blocks}{S12, \S5}{g:self-set}
\sym{$\Theta_1$, $\Theta_2$, $CB$, $IB$}{block sets of the grammar}{S12, \S5}{g:theta-blocks}
\sym{$o$}{output map, model to trees}{S5, Thm.~2}{g:sym-o}
\sym{$\sigma$}{normalizing terms of size $\le M$}{S8, \S3}{g:family}
\sym{$t_k$, $u_k$}{convertible pairs with large joins}{S4, Thm.~4}{g:join-family}
\hline
\multicolumn{4}{@{}l}{\emph{Working labels (Part II)}}\\
\sym{$\approx$}{equal infinitary normal forms; also $\conv$}{LOCK1-structure \S0}{g:wl-bad-pair}
\sym{$B(t,u)$}{bad set of a tree-equal pair}{LOCK1-structure \S0}{g:wl-bad-pair}
\sym{$\kappa(t)$}{kind: 0, 1, 2 or $\bot$}{LOCK1-structure \S0}{g:wl-kind}
\sym{$\mathrm{comp}_\alpha$}{component at address $\alpha$}{LOCK1-CTB \S0}{g:wl-program}
\sym{$rb_i(t)$}{rightmost branch}{LOCK1-CTB \S0}{g:wl-rn}
\sym{$\Kc$, $\Kc_1$}{left-cancellation domains}{LOCK1-CT \S12.6}{g:wl-calK}
\sym{$\Hc$, $\Hc_2$, $\Gc$}{counter grammar of the notes}{LOCK1-CT \S19.2}{g:wl-calH}
\sym{$\Hc^{x}$, $\Hc^{x}_2$}{extended counter grammar}{LOCK1-CT \S20.1}{g:wl-calHx}
\sym{$\Gc^{x}$}{$\Hc^{x}_2\times\Hc^{x}$}{LOCK1-CT \S20.1}{g:wl-calGx}
\sym{$\Gc^{+}$}{pre-entry grammar}{LOCK1-CT \S22.1}{g:wl-calGplus}
\sym{$\Ic$}{inert terms $S\,c_j\,Q$}{LOCK1-CT \S20.1}{g:wl-calI}
\sym{$\mathcal{S}$}{$\Kc$-spined T-stacks (not $\Ter$)}{LOCK1-CT \S25.3}{g:wl-calS}
\sym{$r_1$, $r_2$}{first and second refill}{LOCK1-CT \S15}{g:wl-r1}
\sym{$\mathrm{Ctx}(\zeta)$, $\mathcal X_Z$}{context closure of a self-dead fold}{LOCK1-CT \S26.2}{g:wl-ctx}
\sym{$\mathcal D(Z)$}{heaps built around $Z$ (not $D(v)$, not the 373-model)}{LOCK1-CT \S26.2}{g:wl-cd}
\sym{$\mathrm{Ctx}_\cup$, $\mathcal X$, $\mathcal D^{SD}$}{one-base context class and heaps}{LOCK1-CT \S27.1}{g:wl-cdplus}
\sym{$q(t)$}{state of the normalization automaton NA}{LOCK1-INV \S2}{g:wl-na}
\sym{$\mathcal T$}{tape class of normal pairs}{LOCK1-INV \S I1}{g:wl-tape}
\end{longtable}
\endgroup

\section*{Part I.\texorpdfstring{\quad}{ } Terms of the papers}

\begin{entry}{g:automaton-35}{35-state automaton (\emph{Dance}, Fig.~5.14); $\delta$, $d(t)$, $\circledast$}{2026-11-26}
\src{S1, \S3}
\lean{\texttt{HnfS.}$\delta$, \texttt{HnfS.dance}}
A bottom-up tree automaton of~\citep{dots} with states $Q=\{0,\dots,34\}$, initial state $0$ for $S$,
transitions $\delta$ and one rejecting state $\circledast=34$; it accepts $t$ when $d(t)\ne\circledast$,
where $d(S)=0$ and $d(u\,v)=\delta(d(u),d(v))$. It is the image of the block model under $\varphi$, so
it rejects exactly the terms of block $\bot$ (S1, Cor.~5). \see{g:q10}
\end{entry}

\begin{entry}{g:model-373}{373-element model}{2026-11-26}
\src{S5, Thm.~2}
\lean{\texttt{D2.d2\_prefix}, \texttt{D2.d2\_slaw}, \texttt{D2.d2\_block}}
An $S$-model $D$ with $373$ elements and an output map $o$ such that $\pi'_2(t)=o(\bar t)$ for every
$t\in\Ter$. So the fold determines the depth-two prefix, and the terms with a given depth-two prefix
form a regular tree language. The 38-element block model is a homomorphic image of $D$.
\end{entry}

\begin{entry}{g:sym-A}{$A$}{2026-11-26}
Overloaded.
\sense{S2, \S1} $A=S\,S\,S$ (also S3, Fig.~1); $A\,A\,A$ has no hnf and is not productive.
\sense{S2, Prop.~10} (A), the minimal-pairs hypothesis (\see{g:minimal-pair}).
\sense{S4, \S4} A set $\{S\}$ or $\{S,\,S\,S\}$ in ``$A$-clean'' (\see{g:a-clean}).
\sense{S12, \S2} The stack of a normal pair $(A,W)$ (\see{g:normal-pair}).
\end{entry}

\begin{entry}{g:a-clean}{$A$-clean}{2026-11-26}
\src{S4, \S4}
A property of terms, for $A=\{S\}$ or $\{S,\,S\,S\}$: every subterm $a\,S$ has $a\in A$. Reduction
preserves it, and in an $A$-clean term every step adds at least one leaf.
\end{entry}

\begin{entry}{g:accessible}{accessible}{2026-11-26}
\src{S7, \S2}
\lean{\texttt{HeadRec.acc\_all}, \texttt{HeadRec.bsn}}
A property of labelled terms, defined inductively: a term is accessible if all its
$\prec$-predecessors are accessible, where the predecessors of $b$ are its $N$-step results, its
function part and argument, and $a$ when $b=(l\cdot_i a)\cdot_k r$. S7 uses induction on
accessibility in place of Waldmann's recursive path order. \see{g:labelled-reduction}
\end{entry}

\begin{entry}{g:frakB}{$\Bk$ (the set $\mathfrak B$), pieces, passive arguments}{2026-11-26}
\src{S11, \S5}
\lean{\texttt{S11.InB}, \texttt{S11.RedB}}
The least set containing the pieces $S$, $S\,S$, $Q$, $S\,Q$, $c_j$ and, for passive arguments
$U_1,\dots,U_r\in\Bk$, the families (Fa), (Fx), (FQa), (FQa2), (FSS), (FS) built from $a_m$, $c_j$,
$Q$. Its set of reducts $\mathrm{Red}(\Bk)$ is closed under subterms (S11, Lemmas~16--17), and a
solution of the spine equations in $\Bk$ is $a_k$ (S11, Lemma~18). \see{g:spine-equations}
\end{entry}

\begin{entry}{g:berarducci-tree}{Berarducci tree, root-active, root-stable}{2026-11-26}
\src{S3, \S4}
A tree assigned to a term: $\bot$ if the term is root-active (every reduct reduces to a redex), $S$
if $t\red S$, and otherwise $\mathrm{BeT}(f)\,\mathrm{BeT}(w)$ for a root-stable reduct $f\,w$ (no
reduct of it is a redex). No finite $S$-term is root-active, so the tree has no $\bot$ and equals
$t^{\infty}$~\citep{berarducci}.
\end{entry}

\begin{entry}{g:sym-beta}{$\beta$}{2026-11-26}
Overloaded.
\sense{S1, \S2; S6, \S2; S9, \S4; S12, \S5} The table of the 38-element block model
(\see{g:block-model}). \lean{\texttt{HnfS.}$\beta$}
\sense{S4, \S1} As a subscript, $\conv$ is convertibility (\see{g:conv}).
\sense{S11, \S5} A member of $\Bk$.
\end{entry}

\begin{entry}{g:block-model}{block model (38 elements), block}{2026-11-26}
\src{S1, \S2; reused in S2--S12}
\lean{\texttt{HnfS.blockTable}, \texttt{HnfS.}$\beta$}
A table $\beta$ on the blocks $\{0,\dots,37\}$ that satisfies the $S$-law, with block $0$ for $S$ and
block $36=\bot$ for ``no head normal form''; a term has an hnf iff its block is not $\bot$ (S1,
Thm.~9 with Prop.~4). S6 writes it $\mathbf B$, S9 $B$ (with facts (M1)--(M3), S9, \S4), and S12
$\mathbb M$ (\see{g:m2}). S5 shows that it is a homomorphic image of the 373-element model.
\end{entry}

\begin{entry}{g:bohm-tree}{B\"ohm tree $\mathrm{BT}(t)$}{2026-11-26}
\src{S2, \S2}
A tree assigned to a term: $\mathrm{BT}(t)=\bot$ if $t$ has no hnf, and otherwise
$S\,\mathrm{BT}(a_1)\cdots\mathrm{BT}(a_j)$ for an hnf $S\,a_1\cdots a_j$ of $t$. By S3, Thm.~10 it is
$t^{\infty}$ with every node that has an infinite left spine, and the subtree below it, replaced by
$\bot$; that sentence is not mechanized.
\end{entry}

\begin{entry}{g:sym-bot}{$\bot$}{2026-11-26}
Overloaded.
\sense{S1, \S2} Block $36$, left-absorbing, the block of exactly the terms with no hnf.
\sense{S2, \S2} The root of a B\"ohm tree for a term with no hnf (\see{g:bohm-tree}).
\sense{S3, \S4} The Berarducci tree of a root-active term; it never occurs for finite $S$-terms.
\end{entry}

\begin{entry}{g:bridge}{bridge lemma}{2026-11-26}
\src{S3, Lemma~4}
\lean{\texttt{TopTerm.bridge}}
Lemma. If a reduction from $t$ rewrites the root $k$ times, the head reduction of $t$ does so at
least $k$ times. So top-termination reduces to head reduction.
\end{entry}

\begin{entry}{g:bursts}{bursts (Lemma), $P_k$, $Q_k$}{2026-11-26}
\src{S6, Lemma~7}
\lean{\texttt{BExp.iterS\_run}, \texttt{BExp.fib\_dead}}
Lemma. With $(P_0,Q_0)=(x,y)$ and $(P_{j+1},Q_{j+1})=(Q_j,P_j\,Q_j)$, $S^k u\,x\,y\hred
u\,P_k\,Q_k$ in $k$ head steps. For $k\ge8$ neither $P_k$ nor $Q_k$ has an hnf, and for $k\ge10$ both
are stable frames.
\end{entry}

\begin{entry}{g:sym-c}{$c$, $c_j$}{2026-11-26}
Overloaded.
\sense{S1, \S4} The weight table on pairs (h-block, block) (\see{g:weight}).
\sense{S2, \S3} The class $c(t)$ of a term (\see{g:fold}).
\sense{S4, Thm.~4} The constant in a candidate bound $\max(a,b)+c$.
\sense{S11, \S3} $c_0=G$, $c_{j+1}=Q\,c_j$: the counter pieces of $p_K$ (\see{g:pk-run}).
\lean{\texttt{S11.c}}
\end{entry}

\begin{entry}{g:certificate}{certificate}{2026-11-26}
\sense{S1, \S4} The tables $c$ ($37\times38$) and $\pi$ ($37$ entries), all at most $10$, found by
integer programming and checked by kernel evaluation against inequalities~(3) on all pairs and~(4)
on the $9{,}725$ three-item windows. None exists with every entry at most $9$.
\lean{\texttt{HnfSound.wTable}, \texttt{HnfSound.piTable}}
\sense{S2, \S6} A list of 7,244 claims about symbolic head reduction: 7,118 run claims and 126 split
claims. Once the checker is proved sound (S2, Lemma~12), the list settles hypothesis~(A) (S2,
Prop.~13). \lean{\texttt{HnfCert.checkClaim}, \texttt{HnfCert.claims\_sound}}
\end{entry}

\begin{entry}{g:claim}{claim, symbolic term}{2026-11-26}
\src{S2, \S6}
\lean{\texttt{HnfCert.checkClaim}, \texttt{HnfCert.Inst}}
A symbolic term is a term built from $S$ and variables $X_\sigma$, one per accepted product state;
$u$ is an instance of $X_\sigma$ if $\operatorname{st}(u)=\sigma$. A claim is a symbolic term with a
stratum and a rank, valid if no instance is productive. Split claims case on the two states of an
application; run claims run head reduction symbolically.
\end{entry}

\begin{entry}{g:class-model}{class model, 39-state automaton}{2026-11-26}
\src{S2, \S3}
\lean{\texttt{HnfProd.clsTable}, \texttt{HnfProd.}$\gamma$, \texttt{HnfProd.cls}}
A model $\gamma$ of the $S$-law on the classes $\{0,\dots,38\}$, with class $0$ for $S$ and a
rejecting class $\mathsf r=37$ that absorbs on both sides. A term is accepted iff its class is not
$\mathsf r$, and it is accepted iff it is productive (S2, Thm.~1); the model does not refine the
block model. S3, Cor.~11 reads it as deciding whether $t^{\infty}$ has a node with an infinite left
spine (\texttt{InfNF.nodes\_iff\_cls}).
\end{entry}

\begin{entry}{g:closed-s-terms}{closed $S$-terms $\Ter$, $\Ter^{\infty}$}{2026-11-26}
\src{S1, \S2; S3, \S2}
\lean{\texttt{HnfS.PureS}}
$\Ter$ is the set of closed terms built from $S$ alone by application, written left-associated; it
is the same set throughout the series. S3 writes $\Ter^{\infty}\supseteq\Ter$ for the finite and
infinite binary trees with leaves $S$.
\end{entry}

\begin{entry}{g:closed-walk}{closed walk}{2026-11-26}
\src{S9, \S4 (Lemmas~8 and~9)}
\lean{\texttt{Ncfp.walkB\_eq}, \texttt{Ncfp.comb\_walk}}
A walk of blocks $v_i=b(p\,Y_1\cdots Y_i)$ under right multiplication $v\mapsto\beta(v,w)$ that
returns to its start, as it does when $p\conv p\,Y_1\cdots Y_k$. A rank on the blocks falls along
every arrow except into $\bot$ and the loop $\beta(12,0)=12$, so a closed walk off $\bot$ stays at
block $12$ with every $Y_i=S$.
\end{entry}

\begin{entry}{g:grammar-closure}{closure of the counter grammar}{2026-11-26}
\src{S12, Thm.~18, Cor.~19}
Theorem. A move of $(Z,W)\in\Gc$ with normalizing first tooth is (i) a descent of the stack to
$\mathrm{pop}(Z)$, (ii) for $Z=X$, a move to $(W,S\,W)$, (iii) for $Z=U$, to $(W,D(W))$, or (iv) for
$Z=X$, $W=S\,V$, to $(V,D(S\,V))$, and the new state is in $\Gc$. Moves (ii) and (iii), which copy
the heap into the stack, recur; the notes call the closure $\mathrm{CLOSE}_x$
(\see{g:wl-close-x}).
\end{entry}

\begin{entry}{g:comb}{comb, combs $\sigma_j$, left comb}{2026-11-26}
\sense{S12, \S1} The terms $\sigma_0=S$, $\sigma_{j+1}=S\,\sigma_j$, with $\Cc=\{\sigma_j\}$; their blocks
lie in $CB=\{0,1,2,4,8,17\}$ (S12, Lemma~16). S3 calls them right combs.
\lean{\texttt{Lock1Tables.comb}}
\sense{s8-combs (draft), Def.~1} A sequence of terms in which each step is an hnf step ($q_n$ has hnf
$S\,q_{n+1}$, $S\,a_n\,q_{n+1}$ or $S\,q_{n+1}\,b_n$) or a left step ($q_n$ has no hnf and first stable
frame $q_{n+1}\,w_n$); a left comb has only left steps. S8 does not define ``comb'' alone, and S9,
Thm.~7 is the left-comb case. \lean{\texttt{Ncfp.teeth\_unbounded\_seq}} \see{g:right-comb}
\end{entry}

\begin{entry}{g:comb-context}{comb context}{2026-11-26}
\src{S8, \S2}
A context $C=c_k\,(\cdots(c_1\,[\,])\cdots)$ with $c_i\in\Ter$, $k\ge0$. On a head run, the last item
of every frame is $C[u]$ for a comb context $C$, where $u$ is the term's last item (S8, Lemma~5(b)).
\end{entry}

\begin{entry}{g:conj-inf-conv}{conjecture: $t^\infty$ determines conversion}{2026-11-26}
\src{S6, Conj.~11}
Conjecture. If $t^\infty=u^\infty$ then $t\conv u$. With prefix regularity and the McKinsey--Evans
argument~\citep{mckinsey,evans} it would make convertibility of $S$-terms decidable. The notes call
it Lock~1 (\see{g:wl-lock1}), and CT is one half of it (\see{g:ct}).
\end{entry}

\begin{entry}{g:conj-s7}{conjecture: no conversion with a context of itself}{2026-11-26}
\src{S7, Conj.~6}
Conjecture. No $t\in\Ter$ and context $E\ne[\,]$ satisfy $t\conv E[t]$. It is NCFP (\see{g:ncfp}),
and it would strengthen Statman's theorem.
\end{entry}

\begin{entry}{g:consistent-labelling}{consistent labelling}{2026-11-26}
\src{S7, \S2 (proof of Lemma~4)}
A labelling of $t$ in which every label is at most the right depth of the argument below it.
Labelling each node $f\,x$ by $\mathrm{rd}(x)$ is consistent, and restricted steps lift to $N$-steps
between consistent labellings. \see{g:labelled-reduction}
\end{entry}

\begin{entry}{g:consumption}{consumption}{2026-11-26}
\src{S9, Thm.~1, Cor.~3, Conj.~10}
\lean{\texttt{Ncfp.thm\_NC\_hnf}, \texttt{Ncfp.ncfp\_left\_comb}}
The equation $p\conv p\,Y_1\cdots Y_k$, $k\ge1$, the case $E=[\,]\,Y_1\cdots Y_k$ of NCFP. S9 excludes
it when $p$ has an hnf (Thm.~1) and when the $Y_i$ normalize (Cor.~3); the general case is S9,
Conj.~10. ``Consumption'' is a working name; S9 does not use it.
\end{entry}

\begin{entry}{g:context}{context}{2026-11-26}
\sense{S7, \S4} A term with one hole $[\,]$; $C[u]$ fills it with $u$.
\sense{S8, \S1} A term with any number of holes; $E[p]$ fills every hole with $p$, and $E$ is closed if
it has no hole. \see{g:plain-context}
\end{entry}

\begin{entry}{g:holes-only}{contexts without closed parts}{2026-11-26}
\src{S8, Cor.~17}
\lean{\texttt{Ncfp.ncfp\_holes}}
Corollary. If $E\ne[\,]$ is built from holes by $E'\mapsto S\,E'$ and $E',E''\mapsto S\,E'\,E''$, no
$p$ satisfies $p\conv E[p]$, at any depth.
\end{entry}

\begin{entry}{g:conv}{convertibility $\conv$}{2026-11-26}
\src{S4, \S1}
\lean{\texttt{WordProblem.Conv}, \texttt{WordProblem.conv\_iff\_join}}
The equivalence closure of reduction. The system is orthogonal, so $t\conv u$ iff $t$ and $u$ have a
common reduct (S4, \S1). \see{g:cr}
\end{entry}

\begin{entry}{g:core-form}{core form}{2026-11-26}
\src{S6, \S4}
A presentation of a term as $S^k\,v$ with $v=S$ or $v$ having at least two items; it is unique (by
S6, Lemma~5).
\end{entry}

\begin{entry}{g:counter-grammar}{counter grammar $\Gc$; $\Hc$, $\Ic$, $\Zc$}{2026-11-26}
\src{S12, Def.~1}
A grammar of normal pairs: $\Hc$ and $\Ic$ are the least sets of normal terms with (H0) $X,U\in\Hc$;
(H1) $S\,H\in\Hc$; (H2) $S\,P\,c\in\Hc$ for two-item $P\in\Hc$, $c$ normal; (H3) $S\,P\,Q\in\Hc$ for
one-item $P\in\Hc$, $Q\in\Hc\cup\Ic\cup\{S\}$; (H4) $S\,(S\,P')\,\sigma\in\Hc$ for one-item $P'\in\Hc$,
$\sigma\in\Cc$; (I) $S\,\sigma\,Q\in\Ic$ for $Q\in\Hc$, $\sigma\in\Cc$. $\Zc$ is the set of two-item
members of $\Hc$, and $\Gc=\Zc\times\Hc$; every member of $\Gc$ is a normal pair (S12, Lemma~16). These
are the notes' $\Hc^{x}$, $\Hc^{x}_2$ and $\Gc^{x}$ (\see{g:wl-calHx}), not the notes' $\Hc$.
\end{entry}

\begin{entry}{g:counter-run}{counter run}{2026-11-26}
\src{S12, \S1}
A run of a term $p$ with a state $\rest^e(p)\conv Z\,W$ for some $(Z,W)\in\Gc$ such that every tooth
of $Z\,W$ normalizes. \see{g:ct-counter}
\end{entry}

\begin{entry}{g:cr}{(CR) (Church--Rosser)}{2026-11-26}
\src{S11, \S2; S12, \S1}
\lean{\texttt{church\_rosser}}
Theorem~\citep{terese}. If $t\red u$ and $t\red v$, then $u$ and $v$ have a common reduct. So $\conv$
is the existence of a common reduct, and normal forms are unique. \see{g:normalizes}
\end{entry}

\begin{entry}{g:ct}{CT (converse of the tooth lemma)}{2026-11-26}
\src{S11, Conj.~19; S12, Conj.~22}
Conjecture. If $p,q\in\Ter$ have no hnf and $\tau_n(p)\conv\tau_n(q)$ for all $n\ge1$, then
$p\conv q$. Open; S6's conjecture ($t^\infty=u^\infty$ implies $t\conv u$) would imply it. S11 proves
it when one term is $p_K$ or one of its states, and S12 when one term has a counter run.
\see{g:ct-pk}, \see{g:ct-counter}, \see{g:wl-ct}
\end{entry}

\begin{entry}{g:ct-counter}{CT against counter runs}{2026-11-26}
\src{S12, Thm.~2}
Theorem. Let $(Z,W)\in\Gc$ have only normalizing teeth, and let $p,q$ have no hnf, with
$\rest^e(p)\conv Z\,W$ and $\tau_n(q)\conv\tau_n(p)$ for all $n\ge1$; then $q\conv p$. The proof is
not mechanized, apart from the table facts. \see{g:u-family}, \see{g:x-family}
\end{entry}

\begin{entry}{g:ct-pk}{CT against $p_K$}{2026-11-26}
\src{S11, Thm.~1, Cor.~2}
\lean{\texttt{S11.ct\_pK}, \texttt{S11.ct\_pK\_rest}}
Theorem. Every $q\in\Ter$ without an hnf with $\tau_n(q)\conv\tau_n(p_K)$ for all $n\ge1$ is
convertible with $p_K$; so is every such $q$ against a state $\rest^n(p_K)$, with the teeth shifted
by $n$. The Lean proof uses only the axioms \texttt{propext} and \texttt{Quot.sound}. \see{g:pK}
\end{entry}

\begin{entry}{g:sym-D}{$D$, $d$, $D(v)$}{2026-11-26}
Overloaded.
\sense{S1, \S3} $d(t)$, the state of the 35-state automaton (\see{g:automaton-35}).
\sense{S3, \S2} $d(s,t)=2^{-n}$, the metric on $\Ter^{\infty}$; also a depth bound.
\sense{S5, Thm.~2} $D$, the 373-element $S$-model (\see{g:model-373}).
\sense{S12, \S1} $D(v)=S\,v\,(S\,v)$. \lean{\texttt{Lock1Tables.D}}
\end{entry}

\begin{entry}{g:depth-prefix}{depth-$d$ prefix}{2026-11-26}
\sense{S3, Rem.~8} The prefix of $t^{\infty}$ to depth $d$, with holes below. In Lean it is a relation
$\mathrm{NF}(d,t,P)$, defined by one-level forms of the rules of S3, Prop.~9, and shown to exist, to be
unique and to be unchanged by reduction. \lean{\texttt{InfNF.NF}}
\sense{S5, Def.~1} $\pi'_d(t)$: $\pi'_0=\pi'_{-1}=\bullet$, and for $d\ge1$ one of $S$,
$S\,\pi'_{d-1}(a)$, $S\,\pi'_{d-2}(a)\,\pi'_{d-1}(b)$, $\pi'_{d-1}(f)\,\pi'_{d-1}(w)$, by the shape that
head reduction reaches. It shows the head $S$ of an hnf, which S3's prefix cuts.
\lean{\texttt{D2.Pre}, \texttt{D2.pre\_nf}}
\sense{S6, \S2} $\pi_d$, the prefix of S5 (cited there as ``Def.~2''; S5 numbers it Def.~1).
\lean{\texttt{D2.Pre}, \texttt{Conv.pre\_unique}}
\end{entry}

\begin{entry}{g:sym-E}{$E$}{2026-11-26}
Overloaded.
\sense{S1, \S4} $E=\{0,17\}$, the exempt blocks at which a head step may leave the weight unchanged:
block $0$ is paid by the count of $S$ items, and block $17$ (towers) by $\tau$.
\lean{\texttt{HnfSound.exempt}}
\sense{S7, Conj.~6; S8, \S1} A context (\see{g:context}).
\end{entry}

\begin{entry}{g:conj-entry}{entry conjecture (second refill)}{2026-11-26}
\src{S12, Conj.~23}
Conjecture. In the run of every normal pair whose teeth all normalize, the state after the second
refill is in $\Gc$. With S12, Lemma~7 and Thm.~2 it would give CT against every $p$ whose teeth
normalize; the notes call it R2 (\see{g:wl-r2}).
\end{entry}

\begin{entry}{g:exposed-move}{exposed move}{2026-11-26}
\src{S11, \S1}
Move $n$ of $t$ is exposed if (i) it is of kind (B), (ii) at least one head step reaches the last hnf
$S\,a_k\,b_k$, and (iii) the $y$ with $b_k=y\,a_k$ has no hnf. Then the tooth $y\,a\,w$ reveals the new
state $a\,w$. No move of $p_K$ is exposed. \see{g:moves}
\end{entry}

\begin{entry}{g:f2}{(F2)}{2026-11-26}
\src{S11, \S2; S12, \S1}
\lean{\texttt{Full.run\_total}}
Theorem. The head run of every closed $S$-term reaches an hnf or a stable frame. The cited
proof~\citep{s1,s3} uses the block model and Waldmann's top-termination theorem, and Lean proves it
from mechanized top-termination (\texttt{TopTerm.reach\_all}). S6, \S2 and S9, item (B) state the same
fact.
\end{entry}

\begin{entry}{g:f-step}{f-step, $\red_f$}{2026-11-26}
\src{S7, \S4}
\lean{\texttt{HeadRec.FStep}}
A step whose redex is the whole term or the function part of its parent (not an argument, though
possibly inside one); head steps are f-steps. By S7, Thm.~5 (\texttt{HeadRec.theorem\_H}), if
$t\red_f u$ and $t\red C[u]$, then $C=[\,]$.
\end{entry}

\begin{entry}{g:family}{family, class, type, $\sigma$}{2026-11-26}
\src{S8, \S3}
\lean{\texttt{Ncfp.cor\_X1}, \texttt{Ncfp.ncfp\_X1}}
A collection of classes $Z_i$, each with a type $[Y]$ or $[Y_1,Y_2]$ of classes or $\sigma$ (at least
one a class) that its members' hnf items follow, where $\sigma$ is the set of normalizing terms of
size at most a fixed $M$. S8, Lemma~11 builds a family from $p\conv E[p]$.
\end{entry}

\begin{entry}{g:fin}{$\mathrm{Fin}(t)$}{2026-11-26}
\src{S3, \S3}
\lean{\texttt{TopTerm.Fin}, \texttt{TopTerm.Reach}}
A property of $t$: its head reduction rewrites the root boundedly often. S3 proves it in a positive
form: the head path reaches an hnf or a term whose function part has block $\bot$.
\end{entry}

\begin{entry}{g:finite-index}{finite index theorem, $\simeq_g$}{2026-11-26}
\src{S6, Thm.~9}
\lean{\texttt{Full.theorem\_six} (named after an earlier numbering), \texttt{Full.sig\_congruence}}
Theorem. For a finite model there are relations $\simeq_g$ of finite index on terms (core forms with
equal or periodically equivalent exponents, and cores that agree on values, $\Psi$ and, to grade $g$,
hnf items). If $t\simeq_{77}t'$, then $O(C[t])=O(C[t'])$ for every context $C$.
\end{entry}

\begin{entry}{g:finite-meeting}{finite meeting (L6)}{2026-11-26}
\src{S11, Lemma~6}
\lean{\texttt{S11.lemma6}}
Lemma. If $p,q$ have no hnf, $\rest^n(p)\conv\rest^n(q)$ and $\tau_i(p)\conv\tau_i(q)$ for $i\le n$,
then $p\conv q$. So CT reduces to showing that the states eventually meet.
\end{entry}

\begin{entry}{g:finite-reducts}{finite reduct sets}{2026-11-26}
\src{S4, Lemma~1}
\lean{\texttt{WordProblem.sz\_wstep}, \texttt{WordProblem.mem\_reductsUpTo}}
Lemma. For each $s$, the reducts of $t$ of size at most $s$ form a finite set computable from $t$ and
$s$, since size never decreases along a reduction.
\end{entry}

\begin{entry}{g:fold}{fold, block of a term, class of a term}{2026-11-26}
\sense{S1, \S2} The block of a term: $b(S)=0$, $b(u\,v)=\beta(b(u),b(v))$, a homomorphism into the
block model that is invariant under reduction (S1, Lemma~3). S12, \S5 writes $\fold(t)$, with
$\fold(X)=16$ and $\fold(U)=33$. \lean{\texttt{HnfS.fold}, \texttt{HnfS.fold\_wsteps}}
\sense{S2, \S3} The class of a term: $c(S)=0$, $c(u\,v)=\gamma(c(u),c(v))$, also invariant under
reduction (S2, Lemma~5). \lean{\texttt{HnfProd.cls}, \texttt{HnfProd.wsteps\_cls}}
\sense{S5, \S2; S6, \S2} The map $t\mapsto\bar t$ into an $S$-model, with $S\mapsto s$ and
$t\,u\mapsto\bar t\,\bar u$; convertible terms have equal folds. \lean{\texttt{D2.foldD}}
\end{entry}

\begin{entry}{g:frame}{frame}{2026-11-26}
\sense{S6, \S2; S8, \S1} The presentation $S\,u_1\cdots u_m$ of a term, whose $u_i$ are its items; S8
writes it $[u_1,\dots,u_m]$. \see{g:items}
\sense{S9, \S1; S11, \S1; S12, \S1} A term on the head run of a given term. \see{g:head-run}
\end{entry}

\begin{entry}{g:free-step}{free and non-free head step}{2026-11-26}
\src{S6, \S2}
A head step is free if its first item $u_1$ is $S$ or of the form $S\,v$, and non-free otherwise. The
two bounds on head reduction count non-free steps. \see{g:two-bounds}
\end{entry}

\begin{entry}{g:fuel-adequacy}{fuel adequacy}{2026-11-26}
\src{SIM, Thm.~2}
\lean{\texttt{SelfInterp.fuel\_adequacy\_value}, \texttt{SelfInterp.fuel\_adequacy\_timeout}}
Theorem of~\citep[Thm.~5.2]{davis-si}: $\texttt{whnfF}\,\langle k\rangle\,\langle M\rangle\red
\texttt{just}\,\langle N\rangle$ if $M\Downarrow_n N$ with $n<k$, and $\red\texttt{nothing}$ if
$M\Uparrow$ or $n\ge k$, with $\red$ read as weak reduction. SIM adds the converse and the independence
of the value from the budget (SIM, Cor.~3, \texttt{SelfInterp.value\_independent}).
\end{entry}

\begin{entry}{g:function-part}{function part}{2026-11-26}
\src{S3, \S2}
The part $f$ of an application $f\,w$. Head reduction of a term without an hnf reaches a term whose
function part has block $\bot$ (\see{g:fin}).
\end{entry}

\begin{entry}{g:sym-G}{$G$, $G_e$, $G_1$, $G_2$}{2026-11-26}
Overloaded.
\sense{S6, Lemma~8 (proof)} $G_1$ is the graph on block triples $(\bar x,\bar y,\bar z)$ of root frames
$S\,x\,y\,z$, with an edge to each triple allowed for the next root frame; $G_2$ refines it on a set of
$190$ triples. Ranks on them give the bound $77$ (\see{g:psi-skeleton}).
\sense{S11, Thm.~1} $G=S\,Q\,Q$, also written $c_0$. \lean{\texttt{S11.G}}
\sense{S12, \S7} $G_0=D(S\,Y)$, $G_{e+1}=S\,G_e\,(S\,J_{e+1})$ with $J_{e+1}=S\,(S\,Y)\,(S\,G_e)$: the
heaps of the run of $X\,Y$. Lean's \texttt{Jfam Y e} is $J_{e+1}$, and \texttt{Lock1Tables.G} is an
unrelated list of blocks. \lean{\texttt{Lock1Tables.Gfam}, \texttt{Lock1Tables.Jfam}}
\end{entry}

\begin{entry}{g:sym-gamma}{$\gamma$, $\mathsf r$}{2026-11-26}
\src{S2, \S3}
\lean{\texttt{HnfProd.}$\gamma$}
The table $\gamma$ of the 39-class model and its rejecting class $\mathsf r=37$
(\see{g:class-model}).
\end{entry}

\begin{entry}{g:ground-loop}{ground loop}{2026-11-26}
\src{S7, \S1; S8, \S1}
\lean{\texttt{HeadRec.no\_ground\_loop}}
A reduction $t\to^+u$ in which $t$ occurs in $u$ (possibly $u=t$). There are none in $\Ter$
(\see{g:no-ground-loops}).
\end{entry}

\begin{entry}{g:sym-h}{$h(t)$}{2026-11-26}
\src{S2, \S2}
The unique hnf that $t$ head-reduces to, when there is one; S3, \S2 uses it in the same way.
\see{g:head-step}
\end{entry}

\begin{entry}{g:h-block}{h-block, h-fold}{2026-11-26}
\src{S1, \S2}
\lean{\texttt{HnfSound.AllH}}
A block other than $\bot$; a term is h-fold if its block is one. A walk that ends at an h-block visits
only h-blocks.
\end{entry}

\begin{entry}{g:head-run}{head run}{2026-11-26}
\src{S6, \S2; S8, \S1; S9, \S1; S11, \S1; S12, \S1}
The sequence of head steps from a term; it is deterministic and stops at the first hnf. Its terms are
the frames in the sense of S9, S11 and S12, and every head run reaches an hnf or a stable frame
(\see{g:f2}). It is not the run of a normal pair (\see{g:run-pair}).
\end{entry}

\begin{entry}{g:head-step}{head step, head reduction ($\to_h$, $\hred$)}{2026-11-26}
\sense{S1, \S2} The step that contracts the leftmost redex: if $\operatorname{fr}(t)=[x,y,z]
\mathbin{+\!\!+}\rho$, then $t\to x\,z\,(y\,z)\,\rho$, and the new spine is
$\operatorname{fr}(x)\mathbin{+\!\!+}[z,\,y\,z]\mathbin{+\!\!+}\rho$ (S1, eq.~(1)).
\lean{\texttt{HnfSound.root\_step}}
\sense{S2, \S2; S3, \S2} The same step. S2 writes $\to_h^{*}$ for head reduction, which is
deterministic and stops at an hnf $h(t)$; S3 writes $\to_h$ for one step.
\lean{\texttt{HnfSound.HeadChain}, \texttt{HnfSound.headChain\_uniq}}
\sense{S6, \S2; S8, \S1; S9, \S1; S11, \S1; S12, \S1} For $t=S\,u_1\cdots u_m$ with $m\ge3$,
$t\to u_1\,u_3\,(u_2\,u_3)\,u_4\cdots u_m$, and $\hred$ is its reflexive--transitive closure. S7, \S1
describes it as contracting the redex at the head of the left spine.
\end{entry}

\begin{entry}{g:hnf}{head normal form (hnf)}{2026-11-26}
\src{S1, \S2; used in every later paper}
\lean{\texttt{HnfS.IsHnf}, \texttt{HnfS.HasHnf}}
A closed $S$-term with at most two items: $S$, $S\,a$ or $S\,a\,b$. S1 says that $t$ \emph{has} an hnf
if it reduces to one; S8, S9, S11 and S12 say so if its head run reaches one, and call the first hnf
reached \emph{the} hnf of $t$. The two readings agree (S11, Lemma~3(c)). \see{g:whnf}
\end{entry}

\begin{entry}{g:hole}{hole $\bullet$}{2026-11-26}
\src{S5, Def.~1}
A leaf that marks a cut in a prefix tree; cutting a tree to depth $k$ replaces every node at depth $k$
by $\bullet$. It is not the hole $[\,]$ of a context (\see{g:context}).
\end{entry}

\begin{entry}{g:items}{items}{2026-11-26}
\src{S1, \S2; S6, \S2; S11, \S1; S12, \S1}
The terms $u_1,\dots,u_m$ in the unique presentation $t=S\,u_1\cdots u_m$; $t$ is an hnf iff $m\le2$.
A step is internal if it contracts a redex inside an item (S11, \S2). \see{g:spine}, \see{g:frame}
\end{entry}

\begin{entry}{g:join-family}{join family $t_k$, $u_k$; no additive bound}{2026-11-26}
\src{S4, Thm.~4}
\lean{\texttt{WordProblem.Family.not\_smallJoin\_max\_add}}
Theorem. With $f_0=S\,P$, $f_{k+1}=S\,S\,f_k$, $h_0=S\,Q$, $h_{k+1}=S\,S\,h_k$, $t_k=f_k\,P\,R$ and
$u_k=h_k\,Q\,Q$ (S4's $P$, $Q$, $R$), the convertible terms $t_k$, $u_k$ have every common reduct of
size at least $3k+10$, so no bound $\max(a,b)+c$ works. S4, Conj.~5 says that the smallest common
reduct has size $6k+12$ (checked for $k\le5$).
\end{entry}

\begin{entry}{g:kappa}{$\kappa(x)$}{2026-11-26}
\src{S11, \S3, Lemma~10}
\lean{\texttt{S11.Kappa}}
For $x$ with hnf $S\,u_1\cdots u_m$, $m\ge1$, let $x\cdot1=u_1$; $\kappa(x)$ is the least $k$ with
$x\cdot1^k$ defined and without hnf. It is invariant under $\conv$, and $\kappa(x_n)=i_n+1$. It is not
the block variable $\kappa$ of S12, \S5, nor the kind $\kappa$ of the notes (\see{g:wl-kind}).
\end{entry}

\begin{entry}{g:set-K}{$\Kc$}{2026-11-26}
\src{S12, \S1}
The least set of normal terms that contains every normal term with at most one item and contains
$S\,S\,b$ whenever it contains $b$. For normal $c$ and self-dead normal $W$, $c\,W$ normalizes iff
$c\in\Kc$ (S12, Lemma~11). \see{g:wl-calK}
\end{entry}

\begin{entry}{g:s11-lemmas}{L3--L7}{2026-11-26}
\src{S12, \S1}
S12's names for S11's Lemmas~3 (head runs), 4 (moves), 5 (convertible terms), 6 (finite meeting) and
7 (right suffixes). \see{g:moves}, \see{g:tooth-lemma}, \see{g:finite-meeting}, \see{g:right-suffix}
\end{entry}

\begin{entry}{g:labelled-reduction}{labelled reduction, $N$-step}{2026-11-26}
\src{S7, \S2}
\lean{\texttt{HeadRec.LStep}, \texttt{HeadRec.acc\_all}}
A variant of Waldmann's labelling on terms $S$ and $a\cdot_k b$ with $k\in\mathbb N$. An $N$-step
rewrites $((S\cdot_i x)\cdot_j y)\cdot_k z\to(x\cdot_k z)\cdot_{k+1}(y\cdot_k z)$ when $k<N$. Every
labelled term is accessible, so $N$-steps terminate (S7, Lemma~3).
\end{entry}

\begin{entry}{g:last-frame}{last frame}{2026-11-26}
\src{S8, \S1; S9, \S1}
\lean{\texttt{Ncfp.lemma\_L}}
The frame just before the hnf $h\ne t$ on the head run of $t$. It is always $S\,S\,y\,z$ (written
$[S,y,z]$ in S8), and $h=S\,z\,(y\,z)$ (S8, Lemma~5(a); S9, Lemma~4(a)).
\end{entry}

\begin{entry}{g:leafless}{leafless set, kind, mint, primitive}{2026-11-26}
\src{S8, \S4}
\lean{\texttt{Ncfp.thm\_L}, \texttt{Ncfp.lemma\_PR}}
A set $Z$ of terms each of which has an hnf with one or two items (its kind), all in $Z$. The mint of
a head step from $[x,y,z,\dots]$ is the new item $y\,z$, and a member is primitive if no proper subterm
is in $Z$. By S8, Thm.~16, no nonempty leafless set meets only finitely many $\conv$-classes.
\end{entry}

\begin{entry}{g:left-spine}{left spine}{2026-11-26}
\src{S3, \S2}
\lean{\texttt{InfNF.LeftInf}, \texttt{InfNF.leftInf\_iff}}
The path from a node of an infinite term that follows function edges; it ends at a leaf $S$ or is
infinite. The root of $t^{\infty}$ has an infinite left spine iff $b(t)=\bot$ (S3, Prop.~9), and $t$ is
productive iff no node has one (S3, Thm.~10). It is not the spine $\operatorname{fr}(t)$
(\see{g:spine}).
\end{entry}

\begin{entry}{g:lemma-conv}{Lemma (Conv), ``S6 Lemma~4''}{2026-11-26}
\src{S6, Lemma~4}
\lean{\texttt{Conv.lemma3\_model}, \texttt{Conv.obs\_conv}; for terms \texttt{Ncfp.hnf1\_conv},
\texttt{Ncfp.hnf2\_conv}, \texttt{Ncfp.lemma\_T}}
Lemma. If $t\conv t'$, then $O(t)=O(t')$ over any $S$-model; over terms modulo $\conv$ (S8, Lemma~3;
S9, item (A)), convertible terms have hnfs with equally many items, convertible in order, or first
stable frames with convertible parts. The notes cite it as ``S6 Lemma~4'', the source comments of
s8-combs as ``S6's lemma (Conv)'', and the Lean name carries an older number. \see{g:tooth-lemma}
\end{entry}

\begin{entry}{g:lemma-dstar}{Lemma D$^*$}{2026-11-26}
\src{S8, Lemma~9}
\lean{\texttt{Ncfp.lemma\_Dstar}}
Lemma. If $w=g\,u$ has at least three items, an hnf and $\sz{u}\ge3$, then $g$ has a two-item hnf
$[b,a]$. The hnf of $w$ is then $[u,a\,u]$ if $b=S$, $[a\,u,u\,(a\,u)]$ if $b=S\,S$, and
$[p\,q,q\,(p\,q)]$ otherwise, with $p$ a comb context around $u$.
\end{entry}

\begin{entry}{g:lemma-kinds}{Lemma (kinds)}{2026-11-26}
\src{S6, Lemma~5}
\lean{\texttt{PsiSkel.one\_iff}, \texttt{PsiSkel.one\_item\_of\_fold}}
Lemma. The block of $t$ decides whether $t$ is $S$, of the form $S\,v$, or has at least two items.
\end{entry}

\begin{entry}{g:lemma-ss}{Lemma SS}{2026-11-26}
\sense{S8, Lemma~7} Lemma. If $\varphi$ has $k\ge2$ items and $\varphi\hred[S,y]$, then
$\varphi=[S,c,S,\dots,S]$ and $y=c\,S\cdots S$, so $\sz{\varphi}=\sz{y}+2$. No infinite sequence has
$y_i\hred[S,y_{i+1}]$. \lean{\texttt{Ncfp.lemma\_SS}, \texttt{Ncfp.cor\_SS1}}
\sense{S9, Lemma~5} Lemma. If $\varphi\hred S\,S\,y$, then $\sz{\varphi}=\sz{y}+2$; so there is no
infinite chain $q_n\hred S\,S\,q_{n+1}$. \lean{\texttt{Ncfp.lemma\_SS\_sz}, \texttt{Ncfp.ss\_bound}}
\end{entry}

\begin{entry}{g:literal}{literal}{2026-11-26}
\sense{S8, \S1; S9, \S1; S11, \S1; S12, \S1} A term that is its own hnf, that is, has at most two
items.
\sense{S12, \S2} An entry $i$ of a move of kind (B) that needs no head step to reach $S\,a_i\,b_i$; a
non-literal entry has $b_i=y\,a_i$ for some $y$. \see{g:moves}, \see{g:refill}
\end{entry}

\begin{entry}{g:m2}{(M2)}{2026-11-26}
\src{S12, \S5}
\lean{\texttt{Lock1Tables.hnfless\_iff}}
Fact. A closed $S$-term has no hnf iff its block is $36$ \citep[Prop.~4 and Thm.~9]{s1}, with S11,
Lemma~3(c). S9, \S4 lists the same fact among (M1)--(M3).
\end{entry}

\begin{entry}{g:meet}{meeting of runs}{2026-11-26}
\src{S12, Prop.~20}
Proposition. If $\pi\in\Gc$ and $\pi^\circ$ are normal pairs whose teeth all normalize with the same
normal forms, their runs reach equal states at some index.
\end{entry}

\begin{entry}{g:minimal-pair}{minimal pair, hypothesis (A)}{2026-11-26}
\src{S2, Prop.~10}
\lean{\texttt{HnfQ5.MinimalPairs}, \texttt{HnfQ5.q5\_of\_A}}
Hypothesis (A): for accepted $u,v$ with $c(u\,v)=\mathsf r$ and $b(u\,v)\ne\bot$, the term $u\,v$ is not
productive. It implies that every rejected term is unproductive and depends only on $2{,}757$ pairs
of product states. The certificate proves it (S2, Prop.~13).
\end{entry}

\begin{entry}{g:move}{move $n$}{2026-11-26}
\src{S11, \S1; S12, \S1}
The passage from $\rest^{n-1}(t)$ to $\rest^n(t)$, which produces $\tau_n(t)$; its two kinds are given
by the Moves (\see{g:moves}). In S12 a move is a pop or raising (\see{g:pop}) and may be a refill
(\see{g:refill}).
\end{entry}

\begin{entry}{g:moves}{the Moves (L4)}{2026-11-26}
\src{S11, Lemma~4; S9, Lemma~6; S12, \S2 (as L4)}
\lean{\texttt{S11.lemma4}, \texttt{S11.moves\_A}, \texttt{S11.moves\_B}; for S9 \texttt{Ncfp.lemma\_A},
\texttt{Ncfp.bmove}}
Lemma. Let $f=f^-\,u$ have no hnf. (A) If $f^-$ has no hnf, then $\tau_1(f)=u$ and $\rest(f)=f^-$.
(B) Otherwise $f^-\hred S\,a_0\,b_0$, $w_0=u$, $w_{i+1}=b_i\,w_i$ and $a_i\,w_i\hred
S\,a_{i+1}\,b_{i+1}$ for $i<k$, $a_k\,w_k$ has no hnf, $\tau_1(f)=b_k\,w_k$ and $\rest(f)=a_k\,w_k$.
\see{g:wl-moves}
\end{entry}

\begin{entry}{g:ncfp}{NCFP (no conversion fixed points)}{2026-11-26}
\src{S8 (title, \S6); S7, Conj.~6}
Conjecture. No $p\in\Ter$ satisfies $p\conv E[p]$ for a context $E\ne[\,]$ with a hole. Settled cases:
plain contexts of depth at most $2$ and plain contexts without closed parts (S8), and
$E=[\,]\,x_1\cdots x_k$ with normalizing $x_i$ or with $p$ having an hnf (S9). The notes coined the
acronym NCFP; the papers state the conjecture without it.
\end{entry}

\begin{entry}{g:new-state}{new state, tooth normal form $N$}{2026-11-26}
\src{S12, Lemma~6}
Lemma. For a normal pair with normalizing first tooth the move is of kind (B) with literal entry
$0$. With $A^i=\nf(a_i)$, $B^i=\nf(b_i)$ and $W^i=\nf(w_i)$, the new state is the normal pair
$(A^k,W^k)$, and the tooth has normal form $N=\nf(B^k\,W^k)$. \see{g:normal-pair}
\end{entry}

\begin{entry}{g:no-ground-loops}{No Ground Loops theorem}{2026-11-26}
\src{S7, Thm.~1}
\lean{\texttt{HeadRec.no\_ground\_loop}}
Theorem (Waldmann; constructive in S7). No $t\in\Ter$ has $t\to^+u$ with $t$ a subterm of $u$. By S7,
Cor.~2 (\texttt{HeadRec.no\_head\_recurrence}), no $t$ head-reduces in one or more steps to
$t\,X_1\cdots X_k$; S9 cites the theorem as (C).
\end{entry}

\begin{entry}{g:normal-pair}{normal pair, stack, heap}{2026-11-26}
\src{S12, \S2}
A pair $(A,W)$ of normal terms with $A\,W$ without hnf; then $A=S\,A^0\,B^0$ has two items. $A$ is the
stack and $W$ the heap, and the teeth and moves of $(A,W)$ are those of $A\,W$.
\see{g:wl-normal-pair}
\end{entry}

\begin{entry}{g:normalizes}{normalizes, $\nf(t)$}{2026-11-26}
\src{S8, \S1; S9, \S1; S12, \S1}
A term normalizes if it reduces to a normal form, which is then unique by (CR) and written $\nf(t)$;
every term convertible with it reduces to that form and is no larger (S8, S9). By S12, Lemma~5, every
subterm of a normalizing term normalizes, and the term has an hnf.
\end{entry}

\begin{entry}{g:sym-o}{$o$}{2026-11-26}
\src{S5, Thm.~2}
\lean{\texttt{D2.out2}}
The output map from the 373-element model to trees, with $\pi'_2(t)=o(\bar t)$; S6 writes $o_d$ for
depth $d$.
\end{entry}

\begin{entry}{g:observable}{observable $O(t)$}{2026-11-26}
\src{S6, Def.~3}
\lean{\texttt{Conv.HasObs}, \texttt{Conv.obs\_unique}}
A tuple read off the head run over an $S$-model: $O(t)=(m,\bar u_1,\dots,\bar u_m)$ if the run reaches
the hnf $S\,u_1\cdots u_m$, and $O(t)=(\bot,\bar f,\bar w)$ if it first reaches the stable frame
$f\,w$. \see{g:lemma-conv}
\end{entry}

\begin{entry}{g:sym-P}{$P$, $P_k$, $\mathcal P$}{2026-11-26}
Overloaded.
\sense{S2, Def.~3} $P_k$, the approximants of productivity (\see{g:pk}).
\sense{S3, App.~A} $\mathcal P$, the set of right combs (\see{g:right-comb}); in S3, Prop.~9, $P$ is a
function part of block $\bot$ reached by head reduction.
\sense{S4, \S4} $P=S\,S\,S$.
\sense{S6, Lemma~7} $P_k$, with $Q_k$, the burst pair (\see{g:bursts}).
\end{entry}

\begin{entry}{g:pk}{$P_k$ (approximants of productivity)}{2026-11-26}
\src{S2, Def.~3}
\lean{\texttt{HnfProd.Pk}}
Sets of terms: $P_0=\Ter$, and $t\in P_{k+1}$ when $t$ head-reduces to an hnf whose items all lie in
$P_k$. Equivalently, the B\"ohm tree of $t$ has no $\bot$ above depth $k$; each $P_k$ is invariant
under reduction (S2, Lemma~8).
\end{entry}

\begin{entry}{g:pK}{$p_K$}{2026-11-26}
\src{S11, Thm.~1}
\lean{\texttt{S11.pK}}
The term $S\,(S\,S)\,G\,(Q\,(Q\,G))\,S=Q\,c_0\,c_2\,S$, with $c_0=G$ and $c_{j+1}=Q\,c_j$. It has no
hnf, its states are $\rest^n(p_K)=a_n\,S$ and its teeth $Y_n=b_n\,S$ (S11, Lemma~9), and none of its
moves is exposed. \see{g:pk-run}, \see{g:ct-pk}
\end{entry}

\begin{entry}{g:sym-phi}{$\varphi$}{2026-11-26}
Overloaded.
\sense{S1, eq.~(2)} The map from blocks onto the states of the 35-state automaton, a homomorphism with
$\varphi(a)=\circledast$ iff $a=\bot$. \lean{\texttt{HnfS.}$\varphi$, \texttt{HnfS.phi\_onto}}
\sense{S5, Thm.~2} The homomorphism from the 373-element model onto the blocks.
\lean{\texttt{D2.d2\_block}}
\sense{S8, Lemma~7} A term with at least two items, in Lemma SS (\see{g:lemma-ss}).
\end{entry}

\begin{entry}{g:sym-pi}{$\pi$, $\pi'_d$, $\pi_d$}{2026-11-26}
Overloaded.
\sense{S1, \S4} The potential on h-blocks in the certificate (\see{g:weight}).
\lean{\texttt{HnfSound.piTable}}
\sense{S5, Def.~1; S6, \S2} $\pi'_d(t)$, $\pi_d(t)$: the depth-$d$ prefix with the head $S$ shown
(\see{g:depth-prefix}).
\sense{S12, Prop.~20} $\pi$, a normal pair in $\Gc$ (\see{g:meet}).
\end{entry}

\begin{entry}{g:plain-context}{plain context, depth}{2026-11-26}
\src{S8, \S1}
\lean{\texttt{Ncfp.PlainD}}
A context built from the hole and closed normalizing terms by $E'\mapsto S\,E'$ and
$E',E''\mapsto S\,E'\,E''$; its depth is $0$ for the hole or a closed context and otherwise one more
than the larger depth of $E'$, $E''$. By S8, Thm.~1 (\texttt{Ncfp.ncfp\_depth2}), no $p\conv E[p]$
for a plain $E\ne[\,]$ of depth at most $2$ with a hole.
\end{entry}

\begin{entry}{g:pop}{pop, raising move, $\mathrm{pop}(Z)$}{2026-11-26}
Overloaded.
\sense{S12, \S2} A move of a normal pair is a pop if $k=0$ (it keeps the heap and replaces the stack by
a proper subterm) and raising if $k\ge1$ (the heap grows to a proper right extension).
\sense{S12, Lemma~15} For $Z\in\Zc\setminus\{X,U\}$ with first item $S^m\,P_0$, $P_0\in\Zc$,
$\mathrm{pop}(Z)=P_0$.
\end{entry}

\begin{entry}{g:prefix-regularity}{prefix regularity, $M_d$}{2026-11-26}
\src{S6, Thm.~1}
\lean{\texttt{Tower.prefix\_regularity}}
Theorem. For every $d$ there are a finite $S$-model $M_d$, generated by $s$, and a map $o_d$ to trees
with $\pi_d(t)=o_d(\bar t)$, both computable from $d$. $M_1$ is the block model, and $M_{d+1}$ is
$\Ter$ modulo the syntactic congruence of the observable over $M_d$ (S6, Lemma~10).
\see{g:realizable}
\end{entry}

\begin{entry}{g:product-state}{product state $\operatorname{st}(t)$}{2026-11-26}
\src{S2, \S3}
\lean{\texttt{HnfProd.reachPairs}}
The pair $(c(t),b(t))$ of class and block. Only $115$ pairs occur (the set $R$, S2, Lemma~4(iii)),
$95$ of them accepted, and an accepted class never pairs with block $\bot$.
\end{entry}

\begin{entry}{g:productive}{productive}{2026-11-26}
\sense{S2, Def.~3} A term that lies in every $P_k$ (\see{g:pk}); S2, \S2 shows this equivalent to its
B\"ohm tree having no $\bot$, the definition of~\citep{dots}.
\lean{\texttt{HnfCert.Prod}, \texttt{HnfProd.Pk}}
\sense{S3, \S2} A term whose B\"ohm tree has no $\bot$; equivalently, no node of $t^{\infty}$ has an
infinite left spine (S3, Thm.~10). \lean{\texttt{InfNF.prod\_iff\_nodes}}
\end{entry}

\begin{entry}{g:productive-runs}{productive runs (Lemma), $\varpi$}{2026-11-26}
\src{S6, Lemma~6}
\lean{\texttt{BExp.skeleton\_r1}, \texttt{BExp.r1\_le}}
Lemma. There is a map $\varpi$ from the blocks other than $\bot$ to $\{0,\dots,10\}$ such that the head
run of a term $t$ with an hnf has at most $\varpi(\bar t)$ non-free steps. \see{g:two-bounds}
\end{entry}

\begin{entry}{g:sym-Psi}{$\Psi(t)$}{2026-11-26}
\src{S6, Def.~3}
The pair $(\bar f,\bar w)$ of values of the first stable frame $f\,w$ of a term $t$ without an hnf.
\end{entry}

\begin{entry}{g:psi-skeleton}{$\Psi$-skeleton (Lemma)}{2026-11-26}
\src{S6, Lemma~8}
\lean{\texttt{GCert.psi\_skeleton}, \texttt{PsiSkel.root\_head\_not\_comb}}
Lemma. The head run of a term without an hnf has at most $77$ non-free steps before its first stable
frame. The proof uses ranks on the graphs $G_1$, $G_2$ (\see{g:sym-G}). \see{g:two-bounds}
\end{entry}

\begin{entry}{g:sym-Q}{$Q$, $\mathcal Q$, $Q_k$, $q_n$}{2026-11-26}
Overloaded.
\sense{S1, \S3} $Q=\{0,\dots,34\}$, the states of the 35-state automaton.
\sense{S3, \S3} $\mathcal Q$, the terms that are not right combs; membership is invariant under
reduction. \lean{\texttt{TopTerm.qB}}
\sense{S4, \S4} $Q=S\,(S\,S)$.
\sense{S6, Lemma~7} $Q_k$, with $P_k$, the burst pair (\see{g:bursts}).
\sense{S8, Thm.~2; S9, Thm.~7} $q_n$: the terms of a right comb, resp.\ the chain of first stable
frames.
\sense{S9, \S5} The set of $p$ with an hnf and $p\conv p\,Y_1\cdots Y_k$ for some $k\ge1$; it is empty.
\lean{\texttt{Ncfp.SComb}, \texttt{Ncfp.scomb\_closed}}
\sense{S11, Thm.~1} $Q=S\,(S\,S)$, so $Q\,u\,v\hred S\,(u\,v)\,(v\,(u\,v))$. \lean{\texttt{S11.Q}}
\sense{S12, Def.~1} A metavariable for a member of $\Hc\cup\Ic\cup\{S\}$ in rules (H3) and (I).
\end{entry}

\begin{entry}{g:q5}{Question~5 (productive terms)}{2026-11-26}
\src{S2, Thm.~1}
\lean{\texttt{HnfCert.q5}}
A question of \citet{dots}: how to characterize the productive closed $S$-terms, whether they form a
regular set, and whether productivity is decidable. S2 answers it with a 39-state automaton that
accepts exactly the productive terms (\see{g:class-model}).
\end{entry}

\begin{entry}{g:q10}{Question~10 (head normal forms)}{2026-11-26}
\src{S1, Thm.~1}
\lean{\texttt{HnfSound.q10}}
A question of \citet{dots}: whether their 35-state automaton accepts exactly the closed $S$-terms with
a head normal form. S1 proves that it does: rejection gives no hnf (S1, Cor.~5), and every term whose
block is not $\bot$ reduces to an hnf (S1, Thm.~9, \texttt{HnfSound.sound}).
\end{entry}

\begin{entry}{g:sym-r}{$r$, $R$}{2026-11-26}
Overloaded.
\sense{S1, \S4} $r(t)$, the weight of a term (\see{g:weight}).
\sense{S2, \S3} $\mathsf r=37$, the rejecting class; $R$, the set of 115 reachable product states.
\sense{S4, \S4} $R=S\,S$.
\end{entry}

\begin{entry}{g:realizable}{realizable prefixes}{2026-11-26}
\src{S6, Cor.~2}
\lean{\texttt{Tower.realizable}, \texttt{Tower.realizable\_iff}}
Corollary. For each $d$ and finite tree $P$, the terms with depth-$d$ prefix $P$ form a regular tree
language, and it is decidable whether some closed $S$-term has an infinitary normal form with that
prefix. Lean has the decidability for the prefixes $\pi_d$; the reading as prefixes of $t^\infty$ and
the regular-language claim are not mechanized.
\end{entry}

\begin{entry}{g:refill}{refill}{2026-11-26}
\src{S12, \S8}
A move of a normal pair with at least one non-literal entry. All-literal moves shrink the stack, so
refills recur. \see{g:conj-entry}, \see{g:wl-refill}
\end{entry}

\begin{entry}{g:replay}{replay}{2026-11-26}
\src{S7, \S3}
The repetition of the steps of a loop $t\to t'\red u$ at an occurrence of $t$ inside another term. The
proof of No Ground Loops uses it to contradict accessibility.
\end{entry}

\begin{entry}{g:rest}{rest}{2026-11-26}
\src{S11, \S1; S12, \S1}
\lean{\texttt{S11.rest1}, \texttt{S11.restN}}
The function part $f$ of the first stable frame $f\,w$ of a term $t$ without an hnf: $\rest(t)=f$.
Then $\rest(t)$ again has no hnf, and the terms $\rest^n(t)$ are the states of $t$ (S11).
\see{g:tooth}, \see{g:state}
\end{entry}

\begin{entry}{g:right-comb}{right comb}{2026-11-26}
\sense{S3, Lemma~5} One of the terms $S$, $S\,S$, $S\,(S\,S)$, \dots, whose set S3 writes $\mathcal P$;
a term is a right comb iff its block is in $\{0,1,2,4,8,17\}$, and S1, \S4 calls the members of block
$17$ towers. If $A,B,C$ are not right combs then $b(A\,B\,C)=\bot$; S12 calls these terms combs
$\sigma_j$ (\see{g:comb}). \lean{\texttt{TopTerm.RC}, \texttt{TopTerm.rc\_of\_not\_qt}}
\sense{S8, Thm.~2} A sequence $q_0,q_1,\dots$ in which every $q_n$ has the hnf $S\,q_{n+1}$ or
$S\,a_n\,q_{n+1}$. By S8, Thm.~2, no right comb has teeth of bounded size.
\lean{\texttt{Ncfp.right\_combs}, \texttt{Ncfp.right\_combs\_seq}}
\end{entry}

\begin{entry}{g:right-depth}{right depth $\mathrm{rd}$}{2026-11-26}
\src{S7, \S2}
The length of the right spine: $\mathrm{rd}(S)=0$ and $\mathrm{rd}(f\,x)=\mathrm{rd}(x)+1$. A
restricted step never lowers it.
\end{entry}

\begin{entry}{g:right-suffix}{right suffix, $\sqsubseteq$}{2026-11-26}
\src{S11, \S2, Lemma~7; S12, \S2}
\lean{\texttt{S11.lemma7b}, \texttt{S11.lemma7c}, \texttt{S11.lemma7d}}
A term $e$ such that $t=x_1\,(x_2\,(\cdots(x_d\,e)\cdots))$ for some $d\ge0$, the depth. S12 writes
$e\sqsubseteq t$, and $e\sqsubset t$ if $d\ge1$. S11, Lemma~7 tracks right suffixes under reduction.
\see{g:rs}
\end{entry}

\begin{entry}{g:root-frame}{root frame}{2026-11-26}
\src{S6, \S3}
A frame $S\,x\,y\,z$ with three items and no hnf. The proof of the $\Psi$-skeleton bound tracks the
successive root frames of a head run.
\end{entry}

\begin{entry}{g:root-normal-form}{root normal form}{2026-11-26}
\src{S1, Cor.~10}
Following~\citep{dots}, a closed $S$-term that is $S$, $S\,A$, $S\,A\,B$, or $M_1\,M_2$ with $M_1$
having an infinite head reduction. S1 shows that the last condition is $d(M_1)=\circledast$, so the
property is decidable; only the link to the automaton is mechanized.
\end{entry}

\begin{entry}{g:root-redex}{root redex}{2026-11-26}
\src{S3, \S2}
A closed $S$-term with exactly three items, so that its head step rewrites the root.
\end{entry}

\begin{entry}{g:rs}{(RS)}{2026-11-26}
\src{S12, \S2}
Lemma. If $e\sqsubset t$ and $t$ normalizes, then $e$ normalizes and $\nf(e)\sqsubset\nf(t)$; it comes
from S11, Lemma~7(a). \see{g:right-suffix}
\end{entry}

\begin{entry}{g:run-pair}{run of a normal pair}{2026-11-26}
\src{S12, \S2}
The sequence of states $(A_n,W_n)$ of a normal pair, with tooth normal forms $N_n$;
$\rest^n(A\,W)\conv A_n\,W_n$. \see{g:head-run}
\end{entry}

\begin{entry}{g:pk-run}{run of $p_K$: $c_j$, $a_m$, $b_m$, $i_m$, $x_m$, $Y_n$}{2026-11-26}
\src{S11, \S3, Lemmas~8--9}
\lean{\texttt{S11.am}, \texttt{S11.bm}, \texttt{S11.im}, \texttt{S11.xm}, \texttt{S11.Y},
\texttt{S11.lemma9}}
The spine of $p_K$'s run: $a_0=Q\,c_0\,c_2$, each $a_m$ is $Q\,c_j\,c_n$ or $G\,c_n$, and $a_m\hred
S\,a_{m+1}\,b_{m+1}$ with $b_m=c_{i_m+1}\,a_m$. The index sequence $i_m$ is unbounded, $x_m=c_{i_m}\,a_m$
and $Y_n=b_n\,S$. \see{g:pK}, \see{g:kappa}
\end{entry}

\begin{entry}{g:s-model}{$S$-model, $S$-law}{2026-11-26}
\src{S5, \S2; S6, \S2}
\lean{\texttt{D2.d2\_slaw}}
A set $M$ with a binary operation and an element $s$ that satisfies the $S$-law
$s\,a\,b\,c=a\,c\,(b\,c)$ for all $a,b,c\in M$. S1, Lemma~2(i) states the law for the block table,
and because it holds, the fold is invariant under reduction. \see{g:fold}
\end{entry}

\begin{entry}{g:sym-Sm}{$S^m\,t$}{2026-11-26}
\src{S12, \S1}
The term $S\,(S\,(\cdots(S\,t)))$ with $m$ copies of $S$. S6, Lemma~7 uses $S^k u$ in the same way.
\end{entry}

\begin{entry}{g:scott-encoding}{Scott encoding $\langle t\rangle$}{2026-11-26}
\src{SIM, \S1}
\lean{\texttt{SelfInterp.quote\_inj}}
An encoding of $SKI$ terms that takes four continuations, one per constructor ($S$, $K$, $I$,
application), and applies its own. Encodings of distinct terms have no common reduct (SIM, \S3).
\end{entry}

\begin{entry}{g:self-dead}{self-dead}{2026-11-26}
\src{S12, \S3}
\lean{\texttt{Lock1Tables.self\_dead}}
A closed term $W$ such that $W\,(t\,W)$ has no hnf for every closed $t$. Terms whose block is in
$\Self$ are self-dead, and so is every member of $\Zc$ (S12, Lemma~16(c)). \see{g:self-set},
\see{g:wl-self-dead}
\end{entry}

\begin{entry}{g:self-set}{$\Self$}{2026-11-26}
\src{S12, \S5}
\lean{\texttt{Lock1Tables.SELF}, \texttt{Lock1Tables.mem\_SELF}}
The set of blocks $w$ with $\beta(w,\beta(\kappa,w))=36$ for every block $\kappa$; a term with block
in $\Self$ is self-dead. \see{g:self-dead}
\end{entry}

\begin{entry}{g:size}{size $\sz{t}$}{2026-11-26}
\src{S2, \S2; S4, \S2; S8, \S1; S9, \S1; S11, \S1; S12, \S1}
\lean{\texttt{HeadRec.sz}}
The number of occurrences of $S$ in $t$. A step adds $\sz{z}-1\ge0$ occurrences, so $t\red u$ implies
$\sz{t}\le\sz{u}$ (S4, \texttt{WordProblem.sz\_wstep}; S12, \S1).
\end{entry}

\begin{entry}{g:size-laws}{size laws (F), (L0), (L1), (LA)}{2026-11-26}
\src{S2, Lemma~11}
\lean{\texttt{HnfCertLaws.chain\_factor}, \texttt{HnfCertLaws.L0}, \texttt{HnfCertLaws.L1},
\texttt{HnfCertLaws.LA}}
Lemma. Four laws on how head reduction changes sizes: (F) factors a head reduction; (L0) a head
reduction that ends in a spine of at most one item is empty; (L1) $t\to_h^{*}S\,S\,y$ gives
$\sz{y}=\sz{t}-2$; (LA) $t\to_h^{*}S\,(S\,S)\,y$ gives $\sz{y}\le\sz{t}-2$.
\end{entry}

\begin{entry}{g:join-bound}{Small Join, join bound}{2026-11-26}
\src{S4, Def.~2}
\lean{\texttt{WordProblem.SmallJoin}, \texttt{WordProblem.smallJoin\_decides}}
A property of $g\colon\mathbb N^2\to\mathbb N$: convertible $t,u$ have a common reduct $w$ with
$\sz{w}\le g(\sz{t},\sz{u})$. Classically it holds for some $g$, and if it holds for a computable $g$,
convertibility is decidable (S4, Thm.~3). It is open for computable $g$, and no additive bound works
(\see{g:join-family}).
\end{entry}

\begin{entry}{g:spine}{spine $\operatorname{fr}(t)$}{2026-11-26}
\src{S1, \S2}
\lean{\texttt{HnfSound.spine}}
The list $\operatorname{fr}(t)=[u_1,\dots,u_m]$ of the items of $t=S\,u_1\cdots u_m$. It is neither the
left spine of a node in an infinite tree (\see{g:left-spine}) nor the spine equations of S11.
\end{entry}

\begin{entry}{g:spine-equations}{spine equations, solution from $k$}{2026-11-26}
\src{S11, Def.~14}
\lean{\texttt{S11.SolFrom}}
A term $s_0$ is a solution from $k$ (it satisfies the spine equations of $p_K$ from $k$) if every
$s_j$ has a two-item hnf $S\,s_{j+1}\,t_{j+1}$ with $t_{j+1}\conv b_{k+j+1}$. $a_k$ is one, and within
$\Bk$ the only one (S11, Lemma~18, \texttt{S11.lemma18}). \see{g:frakB}
\end{entry}

\begin{entry}{g:stable-frame}{stable frame}{2026-11-26}
\sense{S6, \S2} A frame $S\,u_1\cdots u_m$ with $m\ge3$ such that $f=S\,u_1\cdots u_{m-1}$ has no hnf,
written $f\,w$ with $w=u_m$; later head steps are head steps of $f$, so $w$ is kept.
\lean{\texttt{Conv.FirstStable}, \texttt{BExp.stable\_unique}}
\sense{S9, \S1; S11, \S1; S12, \S1} A frame $f\,w$ of the head run in which $f$ has no hnf. The two
senses agree in substance, and the first stable frame defines the first tooth and the rest.
\see{g:f2}, \see{g:tooth}
\end{entry}

\begin{entry}{g:state}{state}{2026-11-26}
\sense{S11, \S1} A term $\rest^n(t)$, $n\ge0$, of a term $t$ without an hnf.
\sense{S12, \S2} A normal pair $(A_n,W_n)$ of the run of a normal pair; it stands for a term
convertible with $\rest^n$ of the original term. \see{g:normal-pair}
\end{entry}

\begin{entry}{g:statman}{Statman's theorem}{2026-11-26}
\src{S8, Cor.~18}
\lean{\texttt{Ncfp.no\_fpc}, \texttt{Ncfp.no\_fixed\_point\_combinator}}
Theorem~\citep[Rem.~15]{dots}. $S$ alone defines no fixed-point combinator: no $Y\in\Ter$ has
$Y\,x\conv x\,(Y\,x)$ for all $x$. S8 obtains it as the case $E=S\,[\,]$, $p=Y\,S$, and S7, Conj.~6
would strengthen it.
\end{entry}

\begin{entry}{g:std}{(Std) (standardization)}{2026-11-26}
\src{S11, \S2}
\lean{\texttt{wsteps\_split}}
Theorem~\citep{terese}. If $t\red u$, then $t\hred t'\red u$ with every step from $t'$ to $u$
internal, so $t'$ and $u$ have the same number of items and corresponding items reduce. S12 uses it
through S11.
\end{entry}

\begin{entry}{g:step-adequacy}{step adequacy}{2026-11-26}
\src{SIM, \S3}
\lean{\texttt{SelfInterp.step\_whnf}, \texttt{SelfInterp.step\_step}}
Lemma of~\citep[Lemma~5.1]{davis-si}: the interpreter's \texttt{step} answers \texttt{nothing} on a
weak head normal form and \texttt{just} of the contracted term on a head redex. It holds as stated
under weak reduction, not under head reduction.
\end{entry}

\begin{entry}{g:strong-s7}{strong S7 conjecture}{2026-11-26}
\src{S8, Conj.~19}
Conjecture. There are no $t,u\in\Ter$ and one-hole context $C\ne[\,]$ with $t\red u$ and
$t\red C[u]$. By Church--Rosser it is equivalent to NCFP. \see{g:ncfp}
\end{entry}

\begin{entry}{g:strongly-convergent}{strongly convergent}{2026-11-26}
\src{S3, \S2}
\lean{\texttt{InfSN.not\_weakInf}}
A reduction of length $\omega$ in which, for every $d$, only finitely many steps contract a redex at
depth at most $d$; it has a limit in $\Ter^{\infty}$. Every reduction of length $\omega$ from a finite
$S$-term is strongly convergent (S3, Thm.~6; doubly negated in Lean).
\end{entry}

\begin{entry}{g:sym-T}{$T$, $T_E$}{2026-11-26}
Overloaded; $T$ means $S\,S$ in S2 and S3 but $S\,S\,S$ in S12.
\sense{S2, \S1; S3, Fig.~2 and App.~A} $T=S\,S$; Zachos' term $S\,T\,T\,(S\,T\,T)$ is productive with an
infinite B\"ohm tree. In S2, \S6, $T$, $T_1$, $T_2$ also name symbolic terms.
\sense{S8, \S6 and Fig.~1} $T_E$, the unique finite or infinite tree with $T=E[T]$ for a plain context
$E\ne[\,]$ with normal closed parts; a solution of $p\conv E[p]$ would have it as $p^\infty$.
\sense{S12, \S1} $T=S\,S\,S$; in S11, $T$ names only a frame inside a proof.
\lean{\texttt{Lock1Tables.T}}
\end{entry}

\begin{entry}{g:t-infty}{infinitary normal form $t^{\infty}$}{2026-11-26}
\src{S3, Thm.~7}
\lean{\texttt{InfNF.NF}, \texttt{InfNF.nf\_unique}, \texttt{InfNF.nf\_wsteps}}
The normal form in $\Ter^{\infty}$ reached from $t$ by a strongly convergent reduction; it is shared by
all reducts, finite iff $t$ normalizes, and equal to the Berarducci tree of $t$. S3 proves that the
limit exists, by rounds of outermost steps, and cites its uniqueness~\citep{kksv,terese}. Neither the
existence of the single limit nor its identification with the Lean prefixes is mechanized.
\end{entry}

\begin{entry}{g:t16-t17}{T16, T17 (table facts in S12)}{2026-11-26}
\src{S12, Lemma~14}
\lean{\texttt{Lock1Tables.T16}, \texttt{Lock1Tables.T17}}
Lemma. The block sets $\Theta_1$, $\Theta_2$, $CB$, $IB$ are closed under the grammar rules,
$\beta(\rho,\kappa)=36$ for $\rho\in\Theta_2$ and $\kappa$ a block of a member of $\Hc$ or $\Ic$, and
$\Theta_2\subseteq\Self$; all are checked by kernel evaluation. The labels come from the notes
(\see{g:wl-table-facts}).
\end{entry}

\begin{entry}{g:sym-tau}{$\tau(t)$, $\#_S(t)$}{2026-11-26}
\src{S1, Thm.~9}
\lean{\texttt{HnfSound.hd}, \texttt{HnfSound.nS}}
The tie-breakers of the weight: $\tau(t)$ is the number of leading $S$ applications of the first item
($\tau(S\,v)=1+\tau(v)$), and $\#_S(t)$ the number of items equal to $S$. Neither is the tooth $\tau_n$
of S11 and S12 (\see{g:tooth}).
\end{entry}

\begin{entry}{g:theta-blocks}{$\Theta_1$, $\Theta_2$, $CB$, $IB$}{2026-11-26}
\src{S12, \S5}
\lean{\texttt{Lock1Tables.TH1}, \texttt{Lock1Tables.TH2}, \texttt{Lock1Tables.CB},
\texttt{Lock1Tables.IB}}
The block sets $\Theta_1=\{18,20,27\}$, $\Theta_2=\{16,31,33,34\}$, $CB=\{0,1,2,4,8,17\}$ and
$IB=\{12,22,30\}$: the blocks of the one-item members of $\Hc$, of $\Zc$, of the combs and of $\Ic$.
\end{entry}

\begin{entry}{g:tooth}{tooth, teeth}{2026-11-26}
\sense{S8, Thm.~2} The terms $a_n$ of a right comb, where $q_n$ has the hnf $S\,a_n\,q_{n+1}$.
\sense{s8-combs (draft), Def.~1} The $a_n$, $b_n$ of an hnf comb and the $w_n$ of a left comb.
\sense{S9, Thm.~7} The terms $w_n$, where $q_{n+1}\,w_n$ is the first stable frame on the head run of
$q_n$; their sizes are unbounded. The text does not use the word.
\lean{\texttt{Ncfp.teeth\_unbounded\_seq}}
\sense{S11, \S1; S12, \S1} The terms $\tau_n(t)$: for $t$ without an hnf with first stable frame
$f\,w$, $\tau_1(t)=w$ and $\tau_{n+1}(t)=\tau_1(\rest^n(t))$. S12 often gives a tooth of a normal pair
by its normal form $N$. \lean{\texttt{S11.tooth1}, \texttt{S11.tooth}}
\end{entry}

\begin{entry}{g:tooth-lemma}{tooth lemma (L5)}{2026-11-26}
\src{S11, Lemma~5; S12, \S1 (as L5); cf.\ S6, Lemma~4}
\lean{\texttt{S11.inv\_hnf2}, \texttt{S11.inv\_first}}
Lemma. Let $t\conv t'$. If $t$ has the hnf $S\,u_1\cdots u_m$, then $t'$ has an hnf with the same $m$
and convertible items; if $t$ has no hnf, neither has $t'$, and $\tau_n(t)\conv\tau_n(t')$ and
$\rest^n(t)\conv\rest^n(t')$ for all $n$. Its converse with teeth only is CT, and the notes call part
(b) Lemma~T (\see{g:wl-lemma-t}).
\end{entry}

\begin{entry}{g:top-termination}{top-termination}{2026-11-26}
\src{S3, Thm.~3}
\lean{\texttt{TopTerm.top\_termination}, \texttt{TopTerm.Perpetual}}
Theorem of Waldmann~\citep{waldmann-thesis}. No reduction sequence from a closed $S$-term rewrites the
root infinitely often. S3 mechanizes his proof with the sets of his case analysis read as sets of
blocks; the theorem fails for infinite terms (\see{g:sym-U}).
\end{entry}

\begin{entry}{g:two-bounds}{two bounds on head reduction}{2026-11-26}
\src{S6, \S3 (Lemmas~6 and~8)}
\lean{\texttt{BExp.r1\_le}, \texttt{GCert.psi\_skeleton}}
Lemmas. The head run of a term with an hnf has at most $\varpi(\bar t)\le10$ non-free steps, and the
head run of a term without an hnf has at most $77$ non-free steps before its first stable frame.
Prefix regularity rests on them. \see{g:productive-runs}, \see{g:psi-skeleton}
\end{entry}

\begin{entry}{g:sym-U}{$U$, $U_j$}{2026-11-26}
Overloaded.
\sense{S3, \S3} The infinite term $U=S\,U$; $S\,U\,S\,S$ rewrites the root forever, so top-termination
fails for infinite terms.
\sense{S11, \S5} $U$, $U_i$: passive arguments of members of $\Bk$ (\see{g:frakB}).
\sense{S12, \S1} $U=S\,T\,T$ (block $33$); $U_0=U$, $U_{m+1}=D(U_m)$.
\lean{\texttt{Lock1Tables.U}, \texttt{Lock1Tables.Ufam}, \texttt{Lock1Tables.fold\_U}}
\end{entry}

\begin{entry}{g:u-family}{$U_j\,U_m$ family}{2026-11-26}
\src{S12, Cor.~3}
\lean{\texttt{Lock1Tables.F1}}
Corollary. For all $j,m\ge0$ every tooth of $U_j\,U_m$ normalizes, so CT holds against every $p$
without an hnf that has a state convertible with $U_j\,U_m$. The Lean fact \texttt{F1} gives only that
$U_j\,U_m$ has no hnf.
\end{entry}

\begin{entry}{g:uq}{\texttt{UQ} (unquoter)}{2026-11-26}
\src{SIM, Thm.~4}
\lean{\texttt{SelfInterp.UQ}, \texttt{SelfInterp.UQ\_unquote}}
The 43-atom closed $SKI$ term of~\citep{davis-si} with $\texttt{UQ}\,\langle t\rangle\red t$ for every
$SKI$ term $t$.
\end{entry}

\begin{entry}{g:walk}{walk}{2026-11-26}
\src{S1, \S2}
\lean{\texttt{HnfSound.walk}}
The sequence of blocks $q_0=s$, $q_i=\beta(q_{i-1},b(u_i))$ of a list $[u_1,\dots,u_m]$ from a block
$s$; from $0$ the walk of $\operatorname{fr}(t)$ ends at $b(t)$. \see{g:weight}, \see{g:closed-walk}
\end{entry}

\begin{entry}{g:weak-reduction}{weak reduction}{2026-11-26}
\src{S1, \S2; SIM, \S2}
The compatible closure of $S\,x\,y\,z\to x\,z\,(y\,z)$, written $\to$, with $\red$ its
reflexive--transitive closure. SIM reads the cited paper's $\red$ in the same way, as contraction at
any position, not head reduction. \see{g:head-step}
\end{entry}

\begin{entry}{g:weight}{weight $r(t)$, potential $\pi$}{2026-11-26}
\src{S1, \S4}
\lean{\texttt{HnfSound.wt}, \texttt{HnfSound.wTable}, \texttt{HnfSound.piTable}}
A number read along the walk of the spine from $0$: given weights $c$ on pairs (h-block, block) and
potentials $\pi$ on h-blocks, $r(t)=\sum_i c(q_{i-1},b(u_i))$. When the certificate inequalities hold,
each head step from an h-fold term lowers $(r(t),\#_S(t),\tau(t))$ lexicographically (S1, Thm.~9).
\see{g:certificate}, \see{g:sym-E}
\end{entry}

\begin{entry}{g:whnf}{weak head normal form; $M\Downarrow_n N$, $M\Uparrow$}{2026-11-26}
\src{SIM, \S1}
Notation for closed $SKI$ terms: $M\Downarrow_n N$ means that $M$ reaches the weak head normal form
$N$ in exactly $n$ normal-order contractions, and $M\Uparrow$ that it reaches none.
\end{entry}

\begin{entry}{g:whnfF}{\texttt{whnfF}}{2026-11-26}
\src{SIM, Thm.~1}
\lean{\texttt{SelfInterp.whnfF}, \texttt{SelfInterp.whnfF\_eq\_paper}}
The 618-atom closed $SKI$ term of~\citep{davis-si}: given a fuel numeral $\langle k\rangle$ and an
encoded term, it runs the term to weak head normal form and answers \texttt{just}~$\langle N\rangle$
or \texttt{nothing}. SIM proves that bracket abstraction of the definitions yields the printed term.
\see{g:fuel-adequacy}
\end{entry}

\begin{entry}{g:window}{window lemma}{2026-11-26}
\src{S1, Lemma~6}
\lean{\texttt{HnfSound.window}}
Lemma. A head step changes the spine only in a window of three items: before and after the step, the
walk from $0$ reaches the rest $\rho$ of the spine in the same block, by the $S$-law.
\end{entry}

\begin{entry}{g:word-problem}{word problem for $S$}{2026-11-26}
\src{S4, \S1}
The question whether $\conv$ is decidable on closed $S$-terms, asked by Barendregt in
1975~\citep{barendregt-open}; it is RTA open problem \#97~\citep{rtaloop97}, and Waldmann conjectured
it decidable~\citep[\S5]{waldmann-ic}. S4 reduces it to a computable join bound (\see{g:join-bound}),
and S6 to the conjecture that $t^\infty$ determines conversion (\see{g:conj-inf-conv}).
\end{entry}

\begin{entry}{g:sym-X}{$X$}{2026-11-26}
Overloaded; only S12 fixes $X$ as a term.
\sense{S1, \S2} Generic arguments in the hnf shapes $S\,X$, $S\,X\,Y$.
\sense{S2, \S6} $X_\sigma$, the variable of a symbolic term for any term of accepted product state
$\sigma$ (\see{g:claim}).
\sense{S3, App.~A} $X$, $Y$, $Z$ name the arguments of a root redex $S\,X\,Y\,Z$.
\sense{S7, Cor.~2} $X_1,\dots,X_k$: arbitrary arguments in $t\,X_1\cdots X_k$.
\sense{S8, \S1} The positions in the depth-two shapes $S\,X$, $S\,X\,Y$, each among $p$, $n$, $S\,p$,
$S\,p\,p$, $S\,n\,p$, $S\,p\,n$.
\sense{S12, \S1} $X=S\,T\,S$, of block $16$. \lean{\texttt{Lock1Tables.X}, \texttt{Lock1Tables.fold\_X}}
\end{entry}

\begin{entry}{g:x-family}{$X\,Y$ family}{2026-11-26}
\src{S12, Cor.~4, \S7}
Corollary. For $Y=X$ or $Y=S\,X\,c$ with $c$ normal, every tooth of $X\,Y$ normalizes iff $Y=X$ or
$c\in\Kc$, and then CT holds against every $p$ with a state convertible with $X\,Y$. The run moves in
phases through the heaps $G_e$, $J_{e+1}$ (\see{g:sym-G}).
\end{entry}

\begin{entry}{g:xplus}{Theorem X$^+$}{2026-11-26}
\src{S8, Thm.~10}
\lean{\texttt{Ncfp.thm\_Xplus}}
Theorem. There are no $M\ge1$ and nonempty $Z\subseteq\Ter$ such that every member has a one- or
two-item hnf with items in $Z$ or of size at most $M$, at least one in $Z$, and whose last frame
$[S,y,z]$ (if any) has $\sz{z}\le M$, $y=S$, or $y=S\,c$ with $\sz{c}\le M$.
\end{entry}

\begin{entry}{g:sym-Y}{$Y$, $Y_n$}{2026-11-26}
Overloaded.
\sense{S8, Cor.~18} $Y$, a would-be fixed-point combinator (\see{g:statman}).
\sense{S9, \S1} $Y_1,\dots,Y_k$, the absorbed arguments of $p\conv p\,Y_1\cdots Y_k$
(\see{g:consumption}).
\sense{S11, \S3} $Y_n=b_n\,S=\tau_n(p_K)$. \lean{\texttt{S11.Y}}
\sense{S12, Cor.~4} $Y=X$ or $Y=S\,X\,c$ with $c$ normal, the argument of $X\,Y$.
\end{entry}

\begin{entry}{g:set-Z}{$\Zc$}{2026-11-26}
\src{S12, Def.~1}
The set of two-item members of $\Hc$; their blocks lie in $\Theta_2$. Each $Z\in\Zc$ is self-dead, and
$Z\,t$ has no hnf for $t\in\Hc\cup\Ic$ (S12, Lemma~16). \see{g:counter-grammar}
\end{entry}

\section*{Part II.\texorpdfstring{\quad}{ } Working labels (unpublished notes)}

These labels are working names: a statement may be renamed, repaired or withdrawn, and only the papers
are citable for the mathematics. Where a review repaired a statement, the repaired version is given.
Statuses are those of LOCK1-CT through \S27 (Wizard on \S26: sound; on \S27: sound with gaps, no false
[P]; both 2026-11-26). Repairs W26.1--W26.8 are applied in the note, and W26.9 and W27.1--W27.10 are
pending; the entries use the pending repairs' wording.

\begin{entry}{g:wl-ad}{AD (applicatively dead)}{2026-11-26}
\src{LOCK1-CT \S25.3}
A closed $W$ such that $W\,t$ has no hnf for every closed application $t$; for $W$ with an hnf, $W$ is
AD iff $\fold W\in\{31,34\}$ (Lemma ADF, as repaired, W25.1; proved [P], given (S1), T16, T17 [K],
T20, T22 [C]). Every member of $\Hc^{x}_2$ other than $X$, $U$ is AD. Lemma ADH (over an AD heap a stack
with normalizing first tooth pops or makes the T-refill) is the case of Lemma CDH in which the folds
are those of closed applications (\see{g:wl-cdh}).
\end{entry}

\begin{entry}{g:wl-bad-pair}{bad pair, $\approx$, $B(t,u)$}{2026-11-26}
\src{LOCK1-structure \S0}
A pair $(t,u)$ with $t\approx u$ (that is, $t^\infty=u^\infty$) and $t\not\conv u$: a counterexample
to Lock~1. $B(t,u)$ is the set of descent words $\sigma$ with $t_\sigma\not\conv u_\sigma$, empty or an
infinite leafless subtree (Theorem~1 of the note). Overloaded: in LOCK1-CTB $\approx$ means $\conv$
(\S0, and throughout \S7 by Repair WB2.7). \see{g:wl-kind}
\end{entry}

\begin{entry}{g:wl-cd}{Lemma CD (context-death), $\mathcal D(Z)$}{2026-11-26}
\src{LOCK1-CT \S26.2}
\status{proved [P]; given (S1), T24 [C] (Wizard: sound; gloss as repaired, W26.3)}
Lemma. (1) For $Z$ of self-dead fold, $Z\,t$ has no hnf for every $t\in\mathcal X_Z$, in particular for
every closed $t$ that contains $Z$. (2) $\mathcal D(Z)$ is the least set with $Z\in\mathcal D(Z)$ and
$S\,p\,q\in\mathcal D(Z)$ whenever $p\in\mathcal D(Z)$, $q\in\mathcal D(Z)$, or $p=S\,p'$ with
$p'\in\mathcal D(Z)$; every $H\in\mathcal D(Z)$ contains $Z$, and $H\,t$ has no normal form for
$t\in\mathcal X_Z$ (for example $S\,c\,Z$, $S\,Z\,v$, $S\,(S\,Z)\,v$). $\mathcal D(Z)$ is neither
$D(v)=S\,v\,(S\,v)$ nor the 373-element model $D$; Lemma CD$^{+}$ generalizes it (\see{g:wl-cdplus}).
\end{entry}

\begin{entry}{g:wl-cdplus}{Lemma CD$^{+}$, $\mathcal D^{SD}$, $\mathrm{Ctx}_\cup$, $\mathcal X$}{2026-11-26}
\src{LOCK1-CT \S27.1}
\status{proved [P]; given (S1), T25 [C] (part (1) constructive; as repaired, W27.3)}
Lemma. With $\mathrm{Ctx}_\cup$ the union of the sets $\mathrm{Ctx}(\zeta)$ over $\zeta\in\Self$ (23
folds, not $0$) and $\mathcal X$ the closed terms with fold in $\mathrm{Ctx}_\cup$: $\mathcal X$
contains every closed term with a subterm of fold in $\Self\cup\{36\}$ and is closed under application
on either side, and $s\,t$ has no hnf for $\fold s\in\Self\cup\{36\}$, $t\in\mathcal X$.
$\mathcal D^{SD}$, the least set that contains every closed term of fold in $\Self\cup\{36\}$ and
$S\,p\,q$ whenever $p$ or $q$ is in it or $p=S\,p'$ with $p'$ in it, has $H\,t$ without normal form for
$H\in\mathcal D^{SD}$, $t\in\mathcal X$; it contains every $\mathcal D(Z)$, so Lemma CD is a special case.
\end{entry}

\begin{entry}{g:wl-cdh}{Lemma CDH}{2026-11-26}
\src{LOCK1-CT \S26.2}
\status{proved [P]; given (S1), with label G4 (Wizard: sound)}
Lemma. Over a context-dead heap $W$, a normal pair $(A,W)$ with normalizing first tooth has
$B^0\in\Kc$, and either pops ($A^0$ two-item, $A^0\,W$ without hnf) or has $A^0=T$ and makes the
T-refill to $(W,\nf(B^0\,W))$. It generalizes Lemma ADH (\see{g:wl-ad}).
\end{entry}

\begin{entry}{g:wl-cdh-plus}{Cor.~CDH$^{+}$}{2026-11-26}
\src{LOCK1-CT \S27.1}
\status{proved [P]; given (S1), T25 [C], with label G4}
Corollary. Lemma CDH and Cor.~CDS hold over every heap $W\in\mathcal D^{SD}$, with
$Y=\mathrm{Ctx}_\cup$ (\see{g:wl-cdplus}).
\end{entry}

\begin{entry}{g:wl-cds}{Cor. CDS}{2026-11-26}
\src{LOCK1-CT \S26.2}
\status{proved [P]; given (S1) (Wizard: sound; hypotheses (a) and (c), as repaired, W26.2)}
Corollary. In a normalized run all of whose teeth normalize, if the heap $W_n$ is context-dead, then
$A_n\in\mathcal S$; with depth $r$ and bottom $S\,T\,B$, moves $n+1,\dots,n+r$ are pops and move
$n+r+1$ is the T-refill to $(W_n,\nf(B\,W_n))$. The converse is claimed only for a normal self-dead
$W$.
\end{entry}

\begin{entry}{g:wl-cl1}{CL1, FL1, RF, FF}{2026-11-27}
\src{LOCK1-structure \S5}
\status{open [Conj]; Lock~1 $\Leftrightarrow$ CL1 $\Leftrightarrow$ FL1 and Lock~1 $\Leftrightarrow$ RF $\wedge$ FF proved [P]
(\see{g:wl-theorem-c-cong}, \see{g:wl-ff-block})}
Conjectures. CL1: terms with equal trees in every context are convertible; FL1: terms with equal folds
in every model $M_d$ of S6's tower are convertible; each decides the word problem. RF is residual
finiteness, and FF says that the fold into every finite $S$-model is a function of $t^\infty$.
\end{entry}

\begin{entry}{g:wl-close-x}{$\mathrm{CLOSE}_x$}{2026-11-26}
\src{LOCK1-CT \S19.4 (stated), \S20.1 (proved)}
\status{proved [P], in corrected form (Theorem CNT$^{x}$); given (S1), T1--T4, T10, T16, T17 [K]}
\appears{S12 as Thm.~18 (closure of $\Gc$)}
Theorem. $\Gc^{x}$ is closed under moves with a normalizing tooth. As repaired (W19.4), the invariant
is $W^i\in\Hc^{x}\cup\Ic$ for all rounds $i\le m$, inert only at $i=1$, and $W^m\in\Hc^{x}$.
\see{g:grammar-closure}
\end{entry}

\begin{entry}{g:wl-program}{component sync, program, $\mathrm{comp}_\alpha$}{2026-11-26}
\src{LOCK1-CTB \S0, \S1 (Prop.~VS), \S7.3 (Repair WB2.5)}
\status{Prop.~VS proved [P]}
$\mathrm{comp}_\alpha$ follows an address $\alpha\in\{r,w,1,2\}^*$ through rests, teeth and hnf items and
is a conversion invariant (Fact~C). A program is an address $\alpha$ or a pair $(\alpha,\beta)$, and it
fires at move $n$ if the new state (or $a_k$ and $w_k$) is literally that component of the tooth. If
one program fires at the same index on both sides, the pair is convertible (Prop.~VS).
\end{entry}

\begin{entry}{g:wl-context-dead}{context-dead heap}{2026-11-26}
\src{LOCK1-CT \S26 (result paragraph), \S26.2, \S26.4}
An informal name for a heap $W$ with no normal form against every closed argument in a class closed
under contexts; formally, $W$ meets hypotheses (a) and (c) of Lemma CDH for some set $Y$ of folds. The
main instance is $W\in\mathcal D(Z)$ with $Y=\mathrm{Ctx}(\fold Z)$. The open case of TZX$'$ involves
one-item heaps $S\,w$ with $w\in\mathcal D(Z)\cup\{Z\}$, which \S26.4 calls ``2-context-dead''.
\end{entry}

\begin{entry}{g:wl-ct}{CT (Conj.~CT)}{2026-11-26}
\src{LOCK1-CT header, \S1, \S26.7; LOCK1-structure \S3; OPEN-QUESTIONS (A1)}
\status{open [Conj]; proved [P] against $p_K$, the family of Theorem~F and every run that reaches
$\Gc^{x}$, and [K] against $p_K$ (lean-ski v1.21.0, \texttt{S11.ct\_pK})}
\appears{S11 as Conj.~19; S12 as Conj.~22}
Conjecture. If closed $S$-terms $p,q$ without an hnf have $\tau_n(p)\conv\tau_n(q)$ for every $n\ge1$,
then $p\conv q$; Lock~1 is CT together with NTD (\see{g:wl-lock1-ct-ntd}). For normalizing teeth, CT
follows from the chain TZX$_e$ $\wedge$ $\mathrm{RCL}_2^{r}\Rightarrow\mathrm{R2}_e\Rightarrow
\mathrm{ENTRY}_x\Rightarrow$ CT (Theorem R2C$''$): the first step is proved [P], given (S1), T1--T4,
T10, T12, T16, T17 [K], T20, T21, T24 [C], and the rest [P], given also T5, T15 [K]
(\see{g:wl-r2c2}, \see{g:wl-table-facts}).
\end{entry}

\begin{entry}{g:wl-ct-pk}{CT($p_K$), Cor.~CT$_K$}{2026-11-26}
\src{LOCK1-CTB \S11.3, \S12.5}
\status{proved [P], classical (G4) in the notes, after Wizard and Witchhunter review; mechanized in
lean-ski v1.21.0 with only \texttt{propext} and \texttt{Quot.sound}, deciding hnf by the block model}
\appears{S11 as Thm.~1, with the rests as Cor.~2}
\lean{\texttt{S11.ct\_pK}, \texttt{S11.ct\_pK\_rest}}
Theorem. Every closed $q$ without an hnf with $\tau_n(q)\conv\tau_n(p_K)$ for all $n$ satisfies
$q\conv p_K$. The inputs of LOCK1-CTB \S12.5 are G4, (F2), (CR), (Std), (H), S6 Lemma~4, Lemma~T,
Lemma~L, the Moves, Prop.~1 ($\Leftarrow$) and Lemma~R; the block model enters only through (F2) and
G4 (Repair WB6.5). \see{g:ct-pk}
\end{entry}

\begin{entry}{g:wl-ctx}{$\mathrm{Ctx}(\zeta)$, $\mathcal X_Z$, Lemma CTX}{2026-11-26}
\src{LOCK1-CT \S26.2}
\status{Lemma CTX proved [P]; constructive (Wizard: sound)}
For $\zeta\in\Self$, $\mathrm{Ctx}(\zeta)$ is the least set of reachable folds that contains $\zeta$
and is closed under $y\mapsto\beta(x,y)$ and $y\mapsto\beta(y,x)$ for reachable $x$; it has 20 or 21
elements, 36 among them. Lemma CTX: a closed term with a subterm of fold $\zeta$ has its fold in
$\mathrm{Ctx}(\zeta)$, and the set is closed under application on either side.
$\mathcal X_Z=\{t\text{ closed}:\fold t\in\mathrm{Ctx}(\fold Z)\}$.
\end{entry}

\begin{entry}{g:wl-difference-depth}{difference depth, Lemma HR, Cor.~I}{2026-11-26}
\src{LOCK1-NTD-probe, Round~4}
\status{refuted [C] as monotone quantities; Lemma HR and Cor.~I proved [P]}
A family of measures of where the two head forms of a node pair disagree (address depth, function
depth, argument depth, their sum); none is non-increasing at record turns. Lemma HR: a synchronous
head step keeps the function depth and does not raise the argument depth. Cor.~I: every infinite
descent has infinitely many irregular steps.
\end{entry}

\begin{entry}{g:wl-dxp}{Lemma DX$'$}{2026-11-26}
\src{LOCK1-CT \S26.2}
\status{proved [P]; constructive (Wizard: sound)}
Lemma. Lemma DX of \S12.6 holds with hypothesis (i) weakened to ``$d\,x$ has no normal form for every
$x\in\mathcal X$''.
\end{entry}

\begin{entry}{g:wl-event-e}{(E) (essential loss), Theorem E}{2026-11-26}
\src{LOCK1-NTD-probe \S5.5; LOCK1-INJ \S4.1}
\status{(E) computed [C] on a separated pair and open on bad pairs; Theorem~E proved [P], classical,
as repaired (WI1.1)}
(E) is a synchronous head step after which every image of a minimizing essential leaf becomes
inessential although the pair stays non-convertible. Theorem~E: such a step collapses a new node
$a\,z\conv a'\,z'$ with $a\in\{x,x',y,y'\}$, and one of $a$, $a'$ has a two-item hnf.
\end{entry}

\begin{entry}{g:wl-entry}{ENTRY, $\mathrm{ENTRY}_x$, $\mathrm{HEAP}_x$, TREF}{2026-11-26}
\src{LOCK1-CT \S19.2 (ENTRY), \S19.5 and \S20.4 ($\mathrm{ENTRY}_x$), \S20.2 (Prop.~ENT)}
\status{open [Conj]; $\mathrm{ENTRY}_x\Rightarrow$ CT for normalizing teeth proved [P] (Thm.~CT-G$^{x}$),
given (S1), T1--T5, T10, T12, T15--T17 [K]}
\appears{S12, \S8 (the weaker form after Conj.~23)}
Conjectures. ENTRY: every normalized run all of whose teeth normalize has a state in $\Gc$; 150
complete runs in the pools of \S19 never reach $\Gc$ (counters over a right-comb base) [C].
$\mathrm{ENTRY}_x$ is the same with $\Gc^{x}$, follows from R2 and from $\mathrm{R2}_e$, and by Prop.~ENT (as repaired, W20.6)
holds iff from some $n_0$ on every heap is in $\Hc^{x}$ ($\mathrm{HEAP}_x$) and a later move is a
T-round over a two-item member heap or an XB move over $S\,V$, $V\in\Hc^{x}_2$ (TREF).
\end{entry}

\begin{entry}{g:wl-entry-lin}{Conj.~ENTRY-LIN, entry time $E(n)$}{2026-11-27}
\src{LOCK1-INV \S I7.4, \S I8, \S I9}
\status{open [Conj]; computed [C] for $n\le18$ (431{,}336 pairs for $n\le17$ and 817{,}898 for $n=18$), given (S1), (ALG) and the engine's verdicts}
Conjecture. For a normal pair $(A,W)$ with $A$ two-item, $A\,W$ without hnf and $n=|A|+|W|\ge10$: if every tooth of
its run normalizes, then some state of index at most $\lfloor(n-8)/2\rfloor$ lies in $\Gc^{x}$. On the pools the
family $(h^k(S\,(S\,T)\,(S\,T)),S)$ with $h(x)=S\,x\,S$ attains the bound at $n=10+2k$, and a proof route through a
per-move potential on $n$ is refuted [C]. \see{g:wl-entry} \see{g:wl-calGx}
\end{entry}

\begin{entry}{g:wl-essential}{essential difference, elex, Theorem P5}{2026-11-26}
\src{LOCK1-NTD-probe, Round~5 (\S5.3, \S5.5)}
\status{proved [P]; Theorem P5 as repaired to a containment (WN5.2)}
A disagreement leaf is essential if every position pair above it is non-convertible, and elex is an
ordinal rank on essential leaves. Theorem P5: elex decreases at every extraction on a thin path, so
the residue lies in the thin paths with infinitely many asynchronous reads or essential-loss steps,
plus the perfect pairs. \see{g:wl-event-e}
\end{entry}

\begin{entry}{g:wl-fact-k}{Fact K}{2026-11-26}
\src{LOCK1-structure \S1}
\status{proved [P]}
Fact. $t^\infty$ is $S$, $S\cdot a^\infty$, $S\cdot a^\infty\cdot b^\infty$, or $f^\infty\cdot w^\infty$
with an infinite left spine, according to the kind $\kappa(t)$. Hence $t\approx u$ gives equal descent
domains and $t_\sigma\approx u_\sigma$ for every $\sigma$. It is not Theorem~K (\see{g:wl-theorem-k}).
\end{entry}

\begin{entry}{g:wl-fact-s}{Fact S, Fact S$'$}{2026-11-26}
\src{LOCK1-CTB \S9.2, \S12.3}
\status{proved [P], without the block model; Fact S$'$ scoped to $p_K$'s run (Repair WB6.4)}
\appears{S11 as Lemma~8}
Facts. S: if $a_m\hred S\,a_{m+1}\,b_{m+1}$ for all $m$, then $a_m\,u$ has no hnf and its first stable
frame is $a_{m+1}\,u\,(b_{m+1}\,u)$. S$'$: for closed $U_1,\dots,U_r$, $r\ge1$, the terms $a_m\,\vec U$,
$c_j\,a_m\,\vec U$ and $Q\,a_m\,(a_m\,(Q\,a_m))\,\vec U$ have no hnf; together they replace every
block-model check in the CT($p_K$) chain.
\end{entry}

\begin{entry}{g:wl-ff-block}{FF for $\mathbb B$ (the 38-block fold), Lemma B (towers)}{2026-11-27}
\src{LOCK1-FF (as repaired, WF.1--WF.7)}
\status{proved [K]}
\lean{\texttt{FFBlock.fold\_eqv}, \texttt{FFBlock.eq\_tower}}
Theorem. If $t\approx u$ for closed $S$-terms, then $t$ and $u$ have the same block in S1's 38-block fold.
Lemma~B: a term of block $0$, $1$, $2$, $4$, $8$ or $17$ is $S^kS$ for some $k\ge0$. With Consequence~4 of
LOCK1-CONG, the theorem gives FF for every model $M_d$ of S6's tower, and so Lock~1 $\Leftrightarrow$ CL1
$\Leftrightarrow$ FL1 (\see{g:wl-cl1}); FF for an arbitrary finite $S$-model is not claimed (\see{g:wl-ffall}).
\end{entry}

\begin{entry}{g:wl-ffall}{FF$\forall$, small-model certificates, split and comb rules}{2026-11-27}
\src{LOCK1-FFALL \S T1--T3 (as repaired, WFA.1--WFA.7, WT3.1--WT3.8)}
\status{open [Conj]; 3{,}131 of the 6{,}555 $S$-models of size at most~5 certified [P], given a computed check; 3{,}424 open}
Conjecture. FF$\forall$: for every finite $S$-model $M$, $\mathrm{fold}_M$ is $\approx$-invariant; Lock~1
$\Leftrightarrow$ FF$\forall\wedge\mathrm{RF}$. Each model is tested by a greatest fixpoint on pairs of values,
pruned by the split rule (Lemma~L$'$) and the comb rule (Lemma~C). In 1{,}813 open models an explicit term
realizes the residue on one side along the whole tree (the family $e^jQ\,(e^{n+1}Q)$, $e=S\,S$, $Q=S\,e\,e$),
so no one-sided argument closes them; Conj.~IF (finitely many $i$-steps on a path of hnf nodes) would certify
852 more.
\end{entry}

\begin{entry}{g:wl-calGplus}{$\Gc^{+}$}{2026-11-26}
\src{LOCK1-CT \S22.1}
A grammar of normal pairs that contains $\Gc^{x}$ (Lemma GP0); $T$ is its only atom, and its slots may
hold members of $\Kc$, generalized inert terms $S\,k\,Q$ with $k\in\Kc$, and two-item terms $S\,k\,Q$
with $k,Q\in\Kc$. It is not closed under normalized moves (at least nine leaving moves, W22.5) [C], and
Conj.~START$^{+}$ says that every all-normalizing normal pair lies in it.
\end{entry}

\begin{entry}{g:wl-calGx}{$\Gc^{x}$}{2026-11-26}
\src{LOCK1-CT \S20.1}
\appears{S12 as $\Gc$ in Def.~1; its closure is S12, Thm.~18}
The set $\Hc^{x}_2\times\Hc^{x}$ of normal pairs, closed under moves with a normalizing tooth (Theorem
CNT$^{x}$). A pair one of whose runs reaches $\Gc^{x}$ satisfies CT (Cor.~CT-G$_1$; S12, Thm.~2).
Whether every all-normalizing run enters it is $\mathrm{ENTRY}_x$ (\see{g:wl-entry}).
\end{entry}

\begin{entry}{g:wl-h1}{(H1)}{2026-11-26}
\src{LOCK1-CTB \S7.2}
\status{refuted [P] (Theorem T1 of LOCK1-CTB)}
The frozen-heap case of Theorem KH: from some move on, the heap is one literal term $w$ and every
$\tau_1(Y_n)$ is convertible with it. It never happens, so every run with $p_K$'s teeth is an S-run.
\see{g:wl-t1}
\end{entry}

\begin{entry}{g:wl-calH}{$\Hc$, $\Hc_2$, $\Gc=\Hc_2\times\Hc$ (counter grammar of the notes)}{2026-11-26}
\src{LOCK1-CT \S19.2}
A grammar: $\Hc$ is the least set of normal terms with $X,U\in\Hc$, $S\,H$ for $H\in\Hc$, $S\,P\,c$ for
two-item $P\in\Hc$ and normal $c$, and $S\,P\,Q$ for one-item $P\in\Hc$ and $Q\in\Hc\cup\{S\}$; $\Gc$ is
closed under moves with a normalizing tooth (Theorem CNT, proved [P], given (S1), T1--T4, T10, T16
[K]). It is superseded by $\Hc^{x}$: in the inspected runs over a right-comb base the base persists,
so they would not reach $\Gc$ (inspection through 150 moves, not proved). Overloaded: S12's $\Hc$ and
$\Gc$ are the notes' $\Hc^{x}$ and $\Gc^{x}$ (\see{g:wl-calHx}).
\end{entry}

\begin{entry}{g:wl-calHx}{$\Hc^{x}$, $\Hc^{x}_2$}{2026-11-26}
\src{LOCK1-CT \S20.1 (stated informally in \S19.4)}
\appears{S12 as $\Hc$ and $\Zc$ in Def.~1}
An extension of $\Hc$: rule (H3) admits $Q\in\Hc^{x}\cup\{S\}\cup\Ic$, and (H4) adds $S\,(S\,P')\,Q$ for
one-item $P'\in\Hc^{x}$ and $Q$ a right comb; $\Hc^{x}_2$ is its two-item part. Every member of
$\Hc^{x}_2$ is self-dead and has no hnf when applied to a member (via T16, T17 [K]). Every member other
than $X$, $U$ has fold 31 or 34 (via T20 [C]). \see{g:counter-grammar}
\end{entry}

\begin{entry}{g:wl-heap}{heap, S-run}{2026-11-26}
\src{LOCK1-CTB \S7 (Theorem KH), \S12.4}
The heap of a move is its $w_n$, the last argument carried by the state; S12 uses the word for the
second term of a normal pair (\see{g:normal-pair}). An S-run is a run whose heap is the literal $S$ at
every move from the first B-move on, so its states are $a'_n\,S$ and its teeth $b'_n\,S$.
\end{entry}

\begin{entry}{g:wl-hcx}{Lemma HCX (the heaps stay in $\mathcal X$)}{2026-11-26}
\src{LOCK1-CT \S27.1}
\status{proved [P]; given (S1), constructive}
Lemma. If $(A,W)$ is a normal pair with $W\in\mathcal X$, then every round heap of every move of its
normalized run, every later state heap and every tooth that exists lies in $\mathcal X$. It applies to
$(Z,S\,Z)$ for self-dead $Z$, since $S\,Z$ contains $Z$ (\see{g:wl-cdplus}).
\end{entry}

\begin{entry}{g:wl-ii}{(ii), (ii$'$), the Left part}{2026-11-26}
\src{LOCK1-CTB \S7.2, \S8.2}
\status{proved [P] through (W) (LOCK1-CTB \S11); (ii$'$) $\Leftrightarrow$ (ii) $\Leftrightarrow$
CT($p_K$) after Repair WB3.1}
(ii): if $q$ is an S-run with $p_K$'s teeth, then $a'_n\,S\conv a_n\,S$ for some $n$. (ii$'$): the same
for any spine $a'_{n-1}\hred S\,a'_n\,b'_n$ with $b'_n\conv b_n$ and $a'_n\,S$ without hnf. The
equivalence with (ii) needs a prepended prefix (Repairs WB3.1, WB4.4).
\end{entry}

\begin{entry}{g:wl-calI}{$\Ic$ (inert terms)}{2026-11-26}
\src{LOCK1-CT \S20.1}
\appears{S12 as $\Ic$, rule (I) of Def.~1}
The set $\{S\,c_j\,Q: j\ge0,\ Q\in\Hc^{x}\}$, with $c_j$ the right combs; its members are not in
$\Hc^{x}$, and their folds lie in $IB=\{12,22,30\}$ (T17). Lemma INERT says exactly when an inert term
applied to a member normalizes.
\end{entry}

\begin{entry}{g:wl-kill-test}{kill test}{2026-11-26}
\src{LOCK1-CT \S9.4 (first use) and the kill-test subsections of \S\S18--26}
A computation designed to refute a conjecture: it enumerates starts, often adversarial, runs them
within a budget, and looks for a counterexample. A certified non-normalizing tooth confirms an
eventual-death statement such as TZN for that start, and a complete run is only a candidate against
it (as repaired, W25.4); no kill test proves a statement.
\end{entry}

\begin{entry}{g:wl-kind}{kind $\kappa$, head form, descent word, turn}{2026-11-26}
\src{LOCK1-structure \S0}
The kind $\kappa(t)$ is 0, 1 or 2 (the items of its hnf) or $\bot$ (no hnf), and the head form is the
hnf or the first stable frame $f\,w$. A descent word over $\{i,a,b,r,w\}$ addresses a node $t_\sigma$,
and a turn is any letter other than $r$. Overloaded: in LOCK1-CTB \S8.1, $\kappa$ is the length of a
1-chain, the notion that S11 calls $\kappa(x)$ (\see{g:kappa}).
\end{entry}

\begin{entry}{g:wl-calK}{$\Kc$, $\Kc_1$}{2026-11-26}
\src{LOCK1-CT \S12.6}
\appears{S12, \S1 (in-text definition after Thm.~2)}
$\Kc$ is the least set of normal terms that contains every normal term with at most one item and
contains $S\,S\,b$ whenever it contains $b$; it is exactly the set of normal $c$ with $c\,W$
normalizing over a self-dead normal heap $W$ (LCD(1), LCD(2)). $\Kc_1$, for one-item heaps, adds
$S\,(S\,S)\,b$ and the terms $S\,(S\,(S\,S))\,S$ and $S\,(S\,(S\,S))\,(S\,S)$. \see{g:set-K}
\end{entry}

\begin{entry}{g:wl-lcd1}{LCD(1)}{2026-11-26}
\src{LOCK1-CT \S12.6 (Theorem LCD, item~1); \S22.4}
\status{proved [P]; for concrete heaps given (S1), T10 [K]}
\appears{S12 as Lemma~11(b)}
Lemma. If $W$ is normal and self-dead and $c$ is normal with $c\,W$ normalizing, then $c\in\Kc$. As a
classifier of dying moves (\S22.4, as repaired, W22.8) it explains 36--59\% of the deaths at the move
level, so it is not complete [C].
\end{entry}

\begin{entry}{g:wl-lcd2}{LCD(2)}{2026-11-26}
\src{LOCK1-CT \S12.6 (Theorem LCD, item~2)}
\status{proved [P]}
\appears{S12 as Lemma~11(a)}
Lemma. For every $c\in\Kc$ and every normal $W$, $c\,W$ normalizes: $\nf(c\,W)=c\,W$ if $c$ has at
most one item, and $\nf(S\,S\,b\,W)=S\,W\,\nf(b\,W)$.
\end{entry}

\begin{entry}{g:wl-lcd3}{LCD(3)}{2026-11-26}
\src{LOCK1-CT \S12.6 (Theorem LCD, item~3)}
\status{proved [P]}
Lemma. If $W$ is normal and self-dead, then $c\mapsto\nf(c\,W)$ is injective on $\Kc$: left
cancellation holds over $W$.
\end{entry}

\begin{entry}{g:wl-lcd-one}{$\mathrm{LCD}_1$}{2026-11-26}
\src{LOCK1-CT \S12.6 (Theorem LCD$_1$)}
\status{proved [P]; for concrete heaps given (S1), T11 [K]}
Theorem, the one-item analogue of Theorem LCD and distinct from LCD(1). Over a normal one-item heap,
every $c\in\Kc_1$ makes $c\,W$ normalize with one of five explicit normal forms, and
$c\mapsto\nf(c\,W)$ is injective on $\Kc_1$; over a 2-dead heap, $c\,W$ normalizes only for
$c\in\Kc_1$.
\end{entry}

\begin{entry}{g:wl-lemma-d}{Lemma D}{2026-11-26}
Overloaded: two lemmas in two notes.
\sense{LOCK1-structure \S1} Lemma~D (down) [P]: if $t\conv u$, then $\kappa(t)=\kappa(u)$ and the
corresponding children are convertible. Its converse is Lemma~U (up); LOCK1-INJ uses this sense.
\sense{LOCK1-CTB \S10.2} Lemma~D (twins do not cross classes) [P]: a twin's states never meet the
classes of $p_K$'s states at another index. The final chain does not use it.
\end{entry}

\begin{entry}{g:wl-lemma-p}{Lemma P (persistence)}{2026-11-26}
\src{LOCK1-CTB \S10.3, \S12.4}
\status{proved [P]}
\appears{S11 as Lemma~15}
Lemma. If $t$ contains a reduct of $s_j$ for a solution $s$, then every reduct of $t$ contains a
reduct of some $s_{j'}$ with $j'\ge j$.
\end{entry}

\begin{entry}{g:wl-q2}{Lemma Q2 (literal moves over $S\,w$)}{2026-11-26}
\src{LOCK1-CT \S27.1}
\status{proved [P] for parts (1)--(2); given (S1), T25 [C], with label G4; part (3) is information
(W27.5)}
Lemma. Over a heap $S\,w$ with $w\in\mathcal D^{SD}$, a move whose first three rounds are literal dies,
so a Q move (all rounds before $k$ literal) whose tooth normalizes has $k\le2$, and with $k=2$ it lands
over a heap of $\mathcal D^{SD}$. With $k=1$ the new heap can be $S\,(S\,w)$ or $S\,c\,(S\,w)$, which in
general lie outside $\mathcal D^{SD}$; the note's case list there is incomplete (W27.5).
\end{entry}

\begin{entry}{g:wl-qd}{Lemma QD (descent between refills)}{2026-11-26}
\src{LOCK1-CT \S27.1}
\status{proved [P]; constructive; its content is Cor.~RB1(1) (W27.4)}
Lemma. A pop or a Q move replaces the stack by a proper subterm at least 2 smaller, so from a state
$(A_n,W_n)$ at most $(\sz{A_n}-3)/2$ pops and Q moves come in a row, and the next move is a refill or
dies.
\end{entry}

\begin{entry}{g:wl-rn}{Lemma RN, rightmost branch}{2026-11-26}
\src{LOCK1-CTB \S0, \S4, \S12.2}
\status{proved [P], as repaired (WB1.1, WB1.3)}
\appears{S11 as Lemma~7 (parts (b)--(d) are RN, RN-1, RN-2)}
$rb_i(t)$ is the rightmost branch: the last items of successive head forms. Lemma RN: a right suffix
of $t$ is convertible with a node of the branch before its first normalizing node, or normalizes to a
right suffix of that node's normal form; on an infinite branch it may also sink (WB1.1).
\end{entry}

\begin{entry}{g:wl-lemma-sc}{Lemma SC (S-cancellation)}{2026-11-26}
\src{LOCK1-CTB \S8.2, \S12.4}
\status{proved [P]}
\appears{S11 as Lemma~12}
Lemma. (1) If $t$ has no hnf and $e\,S\conv t\,S$, then $e\conv t$. (2) Under side conditions, from
$e\,S\conv h\,S$ with $h\hred S\,u\,v$ follows $e\hred S\,e_1\,e_2$ with $e_1\,S\conv u\,S$ and
$e_2\,S\conv v\,S$.
\end{entry}

\begin{entry}{g:wl-sol}{Lemma SOL}{2026-11-26}
\src{LOCK1-CTB \S11.3, \S12.4}
\status{proved [P]; index step without the block model (Repair WB6.3)}
\appears{S11 as Lemma~18}
\lean{\texttt{S11.lemma18}}
Lemma. A member of $\Bk$ that is a solution from $k$ is literally $a_k$. \see{g:wl-solution}
\end{entry}

\begin{entry}{g:wl-lemma-t}{Lemma T}{2026-11-26}
\src{LOCK1-structure \S0; LOCK1-CTB \S12.1}
\status{proved [P]}
\appears{S11 as Lemma~5(b)}
Lemma. If $t\conv t'$ have no hnf, then $\tau_1(t)\conv\tau_1(t')$ and $\rest(t)\conv\rest(t')$, so all
teeth and rests agree. Its finite converse is S11, Lemma~6, and CT is the converse with teeth only.
\see{g:tooth-lemma}
\end{entry}

\begin{entry}{g:wl-lock1}{Lock 1}{2026-11-27}
\src{A1-lock1-S3; LOCK1-structure (header); OPEN-QUESTIONS (A1)}
\status{open [Conj]; known when one side normalizes [K]; Lock~1 $\Rightarrow$ the word problem is
decidable proved [P]}
\appears{S6 as Conj.~11 (unnamed); S11, \S6 as ``the conjecture of S6''}
\lean{\texttt{WordProblem.Lock1}, \texttt{WordProblem.lock1\_normalizing}}
Conjecture. For closed $S$-terms, $t^\infty=u^\infty$ implies $t\conv u$; since S3 makes every finite
prefix of $t^\infty$ computable, it alone would make conversion decidable. It holds iff CT and NTD
hold (\see{g:wl-lock1-ct-ntd}); CT for normalizing teeth is reduced via $\mathrm{ENTRY}_x$ and $\mathrm{R2}_e$ to
Conj.~TZX$_e$ and Conj.~$\mathrm{RCL}_2^{r}$, CT for non-normalizing teeth is settled only for $p_K$, and
NTD has evidence only. Lock~1 is also equivalent to CL1, to FL1 and to FF$\forall\wedge\mathrm{RF}$ (\see{g:wl-cl1}, \see{g:wl-ffall}). \see{g:conj-inf-conv}
\end{entry}

\begin{entry}{g:wl-lock2}{Lock 2}{2026-11-26}
\src{A1-lock1-S3 (introduction)}
The other lock of an older two-lock route to the word problem. A1-lock1-S3 records it as closed and
as not needed once the prefixes of $t^\infty$ are computable, which leaves Lock~1; it does not restate
its content.
\end{entry}

\begin{entry}{g:wl-moves}{the Moves (A-move, B-move), frame term}{2026-11-26}
\src{LOCK1-CTB \S0, \S12.1}
\status{proved [P], classical (G4) in the notes; [K] in S11}
\appears{S11 as Lemma~4}
\lean{\texttt{S11.lemma4}}
An A-move and a B-move are the cases (A) and (B) of the Moves (\see{g:moves}). In a B-move, entry $i$
is literal if it needs no head step; otherwise its last frame is $S\,S\,y\,z$, and $y$ is the frame
term of the move.
\end{entry}

\begin{entry}{g:wl-na}{NA (normalization automaton), $q(t)$, Prop.~NA-SINK}{2026-11-26}
\src{LOCK1-INV \S2 (as repaired, WI.1--WI.3)}
\status{conditional [P]; sink side given (ALG) [C] on all 59{,}319 state triples, re-run
independently by the Wizard; accepting side given [Lit]}
A bottom-up tree automaton with the 39 states $0,\dots,38$, where $q(S)=0$,
$q(u\,v)=\delta(q(u),q(v))$ and $34$ is the one absorbing sink; it is isomorphic to the automaton
of~\citep[Fig.~5.11]{dots} (all 1{,}444 entries of its $38\times38$ table checked, the sink printed
as $\infty$). (ALG) is the $S$-law on states, $\delta(\delta(\delta(0,a),b),c)=\delta(\delta(a,c),
\delta(b,c))$. Proposition NA-SINK: if $q(t)=34$ then $t$ has no normal form; if $q(t)\ne34$ then $t$
normalizes, by~\citep[Thm.~9]{dots} and~\citep[Thm.~56]{waldmann-ic}. NA decides normalization of
closed $S$-terms [Lit] (not kernel-checked) and differs from the 39-class productivity model of S2
(\see{g:class-model}).
\end{entry}

\begin{entry}{g:wl-normal-pair}{normal pair, normalized run, rounds}{2026-11-26}
\src{LOCK1-CT \S8, \S9.0}
\appears{S12, \S2}
A pair $(A,W)$ of normal terms with $A\,W$ without hnf; $A$ is the stack and $W$ the heap. If all teeth
normalize, the normalized run $(A_n,W_n,N_n)$ records stack, heap and tooth normal form after move $n$,
with rounds $S\,A^{i+1}\,B^{i+1}=\nf(A^i\,W^i)$ and $W^{i+1}=\nf(B^i\,W^i)$ (Lemma~N0).
\see{g:normal-pair}
\end{entry}

\begin{entry}{g:wl-ntz}{Prop.~NTZ, Cor.~NTZ1}{2026-11-26}
\src{LOCK1-CT \S27.3}
\status{proved [P]; given (S1), T8, T12 [K] (W27.1; the note's label also lists T10, T16, T17 [K],
T20, T22 [C]); Wizard: sound with gaps, run recomputed independently}
Proposition. For $Y=S\,(S\,X)\,(S\,X)$ and $Z=S\,(S\,Y)\,(S\,S)$ (two-item, fold 31, $Z\notin\Hc^{x}\cup
\mathcal S$), every tooth of $(Z,S\,Z)$ normalizes: two Q moves lead to $(X,S\,X\,c)\in\Gc^{x}$ with
$c=S\,S\,(S\,Z)\in\Kc$, and the T-start $(X,Z)$ lands at $(Z,S\,Z)$. Corollary NTZ1: TZX$'$, TZ and
$\mathrm{R1}^{+}$ are false, so the routes to R2 through Theorems R2C and R2C$'$ are void (both remain
true implications); R2 holds on this run.
\end{entry}

\begin{entry}{g:wl-ntd}{NTD (no turning descent)}{2026-11-26}
\src{LOCK1-structure \S3}
\status{open [Conj]; no counterexample in Rounds~1--5 of LOCK1-NTD-probe [C]}
Conjecture. For no $t\approx u$ does $B(t,u)$ contain an infinite path with infinitely many turns.
$\mathrm{NTD}_Z$ (LOCK1-NTD-probe \S5.2) restricts it to descents with tail in $\{i,b,w\}^\omega$.
\see{g:wl-lock1-ct-ntd}
\end{entry}

\begin{entry}{g:wl-pk}{$p_K$ and its counter, ``Theorem K's counter''}{2026-11-26}
\src{LOCK1-CTB \S3, \S12.3}
\appears{S11, \S3}
The term $p_K$ of S11 (\see{g:pK}), of size 24, whose counter index $i_n$ is unbounded and takes each
value finitely often. A run ``runs into Theorem~K's counter'' if some state is literally
$Q\,(Q^jG)\,(Q^mG)\,w$ or $G\,(Q^mG)\,w$ for some heap $w$.
\end{entry}

\begin{entry}{g:wl-per}{Per, B$^+$, LC}{2026-11-26}
\src{LOCK1-CT \S4; LOCK1-CTB \S1 (Remark VS-Per), \S8.1}
\status{open [Conj]; Per $\Leftrightarrow$ CT$_1\wedge$ LC and B$^+\Rightarrow$ Per proved [P] (Prop.~4.1);
all three hold at $p_K$ [P]}
Conjectures, for closed $p$ without an hnf and $k\ge1$. Per: no $p$ has $\tau_{n+k}(p)\conv\tau_n(p)$
for all $n$; B$^+$: no $p$ has its teeth in finitely many classes; LC (the left-comb case of NCFP): no
$p\conv p\,Y_1\cdots Y_k$, equivalently no $\rest^k(p)\conv p$. CT$_1$ is CT on one run.
\end{entry}

\begin{entry}{g:wl-r1}{R1, $r_1$}{2026-11-26}
Overloaded; neither sense is the conjecture $\mathrm{R1}^{+}$ or $\mathrm{R1}_s$.
\sense{LOCK1-CT \S8.1} Lemma~R1 [P]: for closed $b$, $W$ with $W\,(b\,W)$ without hnf, $f=S\,T\,b\,W$ has
no hnf, its move has two rounds with non-literal entry~1, $\tau_1(f)=S\,W\,(b\,W)$ and
$\rest(f)=W\,(b\,W)$. ``R1 moves'' and ``R1 loads'' refer to it.
\sense{LOCK1-CT \S15} $r_1$: the index of a run's first refill; a move is visible if its index exceeds
$r_1$, and $r_2$ (also $m_2$) is the second refill.
\end{entry}

\begin{entry}{g:wl-r1-plus}{$\mathrm{R1}^{+}$}{2026-11-26}
\src{LOCK1-CT \S23.3, \S27.3 (Cor.~NTZ1)}
\status{refuted [P] (Cor.~NTZ1); given (S1), T8, T12 [K] (W27.1)}
Conjecture, now false. In a normalized run all of whose teeth normalize, the state right after every
refill that is not of class SSc is in $\Gc^{x}$; the first refill of the T-start $(X,Z)$ of Prop.~NTZ
lands outside $\Gc^{x}$. Theorem R2C (TZ $\wedge$ VB $\wedge$ RCL $\Rightarrow\mathrm{R1}^{+}\Rightarrow$
R2) stays a true implication with a false hypothesis (\see{g:wl-ntz}).
\end{entry}

\begin{entry}{g:wl-r1-s}{$\mathrm{R1}_s$}{2026-11-26}
\src{LOCK1-CT \S23.2 (Cor.~PREF1)}
\status{open [Conj]; R2 $\Rightarrow\mathrm{R1}_s$ proved [P]}
Conjecture. For every normal pair $(Y,H)$ with $H=\nf(b\,S)$ and all teeth normalizing, the state
right after its first refill is in $\Gc^{x}$. A prefix puts one SSc refill in front of $(Y,H)$
(Prop.~PREF), so R2 implies $\mathrm{R1}_s$.
\end{entry}

\begin{entry}{g:wl-r2}{R2}{2026-11-26}
\src{LOCK1-CT \S22.3, \S26.7, \S27.3--27.6}
\status{open [Conj]; not refuted, and it holds on the run of Prop.~NTZ [P]; no known sufficient set}
\appears{S12 as Conj.~23}
Conjecture. In a normalized run all of whose teeth normalize, with $m_2$ the index of the second refill,
$(A_{m_2},W_{m_2})\in\Gc^{x}$; R2 implies $\mathrm{ENTRY}_x$ (Prop.~R2E). The routes through Theorems R2C
and R2C$'$ are void, since TZ is false (Cor.~NTZ1); given TZX$_e$, R2 holds for runs whose first refill
is T-class or XB (Cor.~R2T), and the weaker $\mathrm{R2}_e$ takes its place in the chain to CT
(\see{g:wl-r2e}, \see{g:conj-entry}).
\end{entry}

\begin{entry}{g:wl-r2c2}{Theorem R2C$''$, Cor.~R2T}{2026-11-26}
\src{LOCK1-CT \S27.4}
\status{proved [P]; given (S1), T1--T4, T10, T12, T16, T17 [K], T20, T21, T24 [C]; the step to CT also
needs T5, T15 [K]; Wizard: verified}
Theorem. TZX$_e$ $\wedge$ $\mathrm{RCL}_2^{r}\Rightarrow\mathrm{R2}_e$, and $\mathrm{R2}_e\Rightarrow
\mathrm{ENTRY}_x\Rightarrow$ CT for normalizing teeth. Corollary R2T: given TZX$_e$, R2 holds for every
such run whose first refill is T-class or XB, and then $\mathrm{RCL}_2$ holds for it when entry comes
before the second refill.
\end{entry}

\begin{entry}{g:wl-r2e}{$\mathrm{R2}_e$}{2026-11-26}
\src{LOCK1-CT \S27.4}
\status{open [Conj]; no candidate in the 253,214 normal pairs of \S27.5 [C]}
Conjecture. In a normalized run all of whose teeth normalize, some state with index at most $m_3$, the
index right after the third refill, is in $\Gc^{x}$. R2 implies it and it implies $\mathrm{ENTRY}_x$;
it is not shown to imply R2 (W27.6). \see{g:wl-r2c2}
\end{entry}

\begin{entry}{g:wl-rcl}{RCL, $\mathrm{RCL}_2$}{2026-11-26}
\src{LOCK1-CT \S23.1 (RCL), \S24.2 ($\mathrm{RCL}_2$)}
\status{open [Conj]; the case-(a) half of RCL after $r_1$ proved [P] (Cor.~DP1)}
Conjectures. RCL: in a normalized run all of whose teeth normalize, every refill is of class T, SSc or
XB. $\mathrm{RCL}_2$: in such a run, a case-(b) second refill is the move of Lemma~XB ($X$ over $S\,V$
with $V$ two-item and $V\,D(S\,V)$ without hnf); the census class XB is broader than Lemma~XB (W23.9).
\see{g:wl-refill}, \see{g:wl-rcl2r}
\end{entry}

\begin{entry}{g:wl-rcl2r}{$\mathrm{RCL}_2^{r}$}{2026-11-26}
\src{LOCK1-CT \S27.4, \S27.5}
\status{open [Conj]; confirmed only vacuously [C], with 136 undecided budget stops in its scope (W27.9)}
Conjecture. In a normalized run all of whose teeth normalize and whose first refill is neither T-class
nor the move of Lemma~XB, a case-(b) second refill is the move of Lemma~XB. In the census of 253,214
normal pairs no case-(b) second refill follows such a first refill, and 136 runs with an SSc first
refill stop at the budget undecided. \see{g:wl-r2c2}
\end{entry}

\begin{entry}{g:wl-rec}{REC}{2026-11-26}
\src{LOCK1-CT \S11.5}
\status{open [Conj]; no counterexample in the pools of \S11 [C]}
Conjecture. A normal pair all of whose teeth normalize makes case-(a) refills at infinitely many moves.
In every pool of the notes a case-(b) refill is never followed by another (Conj.~NBB).
\end{entry}

\begin{entry}{g:wl-rec-k}{REC-K}{2026-11-26}
\src{LOCK1-CTB \S7.3 (Repair WB2.5)}
\status{open [Conj]; no counterexample in the small pools of LOCK1-CTB \S7.3 [C]}
Conjecture. In every run with infinitely many non-normalizing teeth that does not run into Theorem~K's
counter, a program fires infinitely often. \see{g:wl-program}
\end{entry}

\begin{entry}{g:wl-rec-t}{$\mathrm{REC}_T$}{2026-11-26}
\src{LOCK1-CT \S17.3}
\status{open [Conj]; $\mathrm{REC}_T\Rightarrow$ CT for normalizing teeth proved [P] (Thm.~CT-T), given
(S1), T1--T5, T10, T12, T15 [K]}
Conjecture. For every $n_0$ there is $n>n_0$ such that move $n$ is a TEE-shaped case-(a) refill (new
heap $S\,A_p$ or $S\,A_p\,(S\,A_p)$), as repaired (W17.3). It holds on the family of Theorem~F and along
every run that reaches $\Gc$ or $\Gc^{x}$.
\end{entry}

\begin{entry}{g:wl-rec-t-ntd}{REC-T-NTD, record turn}{2026-11-26}
\src{LOCK1-NTD-probe, Round~3}
\status{open [Conj]; no counterexample in Round~3 [C]}
Conjecture. In every descent through a bad set, only finitely many turns land on a new record of
$\min(\sz{t_\sigma},\sz{u_\sigma})$ (record turns). Since an infinite descent has infinitely many
records, it implies NTD.
\end{entry}

\begin{entry}{g:wl-refill}{refill; case-(a), case-(b) refill; classes}{2026-11-26}
\src{LOCK1-CT \S9.0, \S9.1, \S10.3, \S11.1, \S23.1}
\appears{S12, \S8 (refill)}
A move of a normal pair that is not all-literal; with $i$ its last non-literal entry, it is of case
(a) if $i=k$ (then the tooth is $N=S\,A'\,W'$, Thm.~TR) and of case (b) if $i<k$, and refills recur
(Cor.~RB1). By the last-round stack $S\,p\,q$, case-(a) refills fall into the classes T ($p=q=S$), SSc
($p=S\,S$, heap $S$; first refill only) and a-other, and case-(b) refills into XB (Lemma~XB: $X$ over a
one-item heap) and b-other. \see{g:refill}
\end{entry}

\begin{entry}{g:wl-calS}{$\mathcal{S}$ ($\Kc$-spined T-stacks)}{2026-11-26}
\src{LOCK1-CT \S25.3, \S26.2}
The least set with $S\,T\,B\in\mathcal{S}$ for $B\in\Kc$ and $S\,Y\,B'\in\mathcal{S}$ for
$Y\in\mathcal{S}$, $B'\in\Kc$; the bottom of $Y$ is the $S\,T\,B$ reached by $r\ge0$ first-argument
steps, and $r$ is its depth. Over an AD heap (Cor.~ADS), and more generally over a context-dead heap
(Cor.~CDS), a stack of an all-normalizing run lies in $\mathcal{S}$, pops $r$ times and makes the
T-refill. Overloaded: it is not the set $\Ter$ of closed $S$-terms, also written $\mathcal S$.
\end{entry}

\begin{entry}{g:wl-s1}{(S1), G4, (F2) as inputs}{2026-11-26}
\src{LOCK1-CTB \S0, \S11.8; LOCK1-CT header}
Inputs named in statuses. ``(S1)'' or ``given S1'' means S1's 38-block model with ``no hnf
$\Leftrightarrow$ block 36'', and G4 marks the use of excluded middle on ``has an hnf'', which that
model decides. (F2) is as in the papers (\see{g:f2}). \see{g:block-model}
\end{entry}

\begin{entry}{g:wl-self-dead}{self-dead, 2-dead}{2026-11-26}
\src{LOCK1-CT \S12.6}
\appears{S12, \S3 (self-dead)}
A term $W$ is self-dead if $W\,(t\,W)$ has no hnf for every closed $t$ (\see{g:self-dead}), and 2-dead
if $W\,p\,q$ has no hnf for all closed $p$, $q$. Their blocks are given by table facts T10 and T11, and
every stack of a normal pair is self-dead (Lemma~SD, T12).
\end{entry}

\begin{entry}{g:wl-solution}{solution from $k$}{2026-11-26}
\src{LOCK1-CTB \S10, \S12.4}
\appears{S11 as Def.~14}
\lean{\texttt{S11.SolFrom}}
The notes' name for a solution of the spine equations (\see{g:spine-equations}). It is invariant under
conversion, and $a'_k$ is one for $k$ past the first B-move of any S-run with $p_K$'s teeth.
\end{entry}

\begin{entry}{g:wl-spine-pair}{spine pair, twin}{2026-11-26}
\src{LOCK1-structure \S3; LOCK1-CTB \S\S8--10}
A counterexample to CT: $p$, $q$ without hnf with all teeth pairwise convertible but $\rest^n(p)
\not\conv\rest^n(q)$ for every $n$; its bad set is the path $r,rr,rrr,\dots$. A twin of $p_K$ is a spine
pair with one side $p_K$, and by CT($p_K$) there is none.
\end{entry}

\begin{entry}{g:wl-status-tags}{status tags [P], [C], [K], [Conj], [Sketch], [Lit]}{2026-11-26}
\src{LOCK1-CT header (as repaired, W9.4, W13.8) and \S24; LOCK1-INV status log, WI.2 ([Lit])}
Tags of the notes. [P]: proved in the note, unmechanized, with the review state recorded per section
in the note's status log; ``[P], given $X$'' is conditional on the inputs $X$, typically (S1) and table
facts. [C]: computed, with command and log, and able only to refute; [K]: kernel-checked in Lean on
lean-ski master; [Conj]: conjectured; [Sketch]: an argument not checked step by step; [Lit]: a published theorem, cited and not re-proved. ``Refuted [P]''
means shown false by a proof, and ``refuted [C]'' by a computed counterexample.
\end{entry}

\begin{entry}{g:wl-sync}{Sync$\bot$, exposed B-move}{2026-11-26}
\src{LOCK1-CTB \S1}
\status{Prop.~S$\bot$ proved [P]}
\appears{S11, \S1 (exposed move)}
A B-move is exposed if its last entry is non-literal and its frame term has no hnf
(\see{g:exposed-move}). Sync$\bot$ (Prop.~S$\bot$): if the first $n$ teeth agree and move $n$ is
exposed on both sides, the pair is convertible. Theorem~K shows that exposure need not recur.
\end{entry}

\begin{entry}{g:wl-t-refill}{T-refill, T-round, T-start, T-stack}{2026-11-26}
\src{LOCK1-CT \S20.2, \S23.1, \S24.2, \S25.3}
A T-round is a round $i$ of a move with $A^i=T=S\,S\,S$; over a self-dead $W^i$ it ends the move as a
case-(a) refill to $(W^i,\nf(B^i\,W^i))$, the heap becoming the stack (Lemma~TRND). A T-refill is a
case-(a) refill whose last-round stack is $T$, and every case-(a) refill after the first is one
(Lemma~DP, Cor.~DP1, proved [P]). A T-start is a pair $(S\,T\,b,Z)$, and a T-stack a stack $S\,T\,B$.
\end{entry}

\begin{entry}{g:wl-t1}{T1}{2026-11-26}
Overloaded: a theorem of LOCK1-CTB and a table fact of LOCK1-CT.
\sense{LOCK1-CTB \S8.1} Theorem~T1 [P; without the block model after WB6.5]: if
$\tau_1(Y_n)\conv\tau_1(Y_{n'})$ then $i_n=i_{n'}$, the counter index being read off a 1-chain (Lemmas
T1a, T1b). Hence (H1) is false; S11 has it as Lemma~10.
\sense{LOCK1-CT \S10} Table fact T1 [K] (\see{g:wl-table-facts}). \lean{\texttt{Lock1Tables.T1}}
\end{entry}

\begin{entry}{g:wl-table-facts}{table facts T1--T25}{2026-11-26}
\src{LOCK1-CT \S10.0 (T1--T5), \S11.4 (T7), \S12.0 (T8--T11), \S13.2 (T12), \S14.1 (T13), \S16.2
(information line, unnamed), \S17.2 (T15), \S19.2 (T16), \S20.1 (T17), \S21.4 (T18), \S22.4 (T19),
\S23.5 (T20), \S24.1 (T21), \S25.3 (T22), \S26.1 (T23), \S26.2 (T24), \S27.1 (T25)}
\lean{\texttt{Lock1Tables.T1}, \dots, \texttt{Lock1Tables.T17}}
\appears{S12 as Lemma~14 (T16, T17)}
Finite checks on S1's 38-block table, which with (S1) become statements about terms. The [K] items are
in \texttt{Ski/Lock1Tables.lean}: T1--T13 and T15 since lean-ski v1.18.0, T16 since v1.19.0, T17 since
v1.20.0; there is no T6 and no T14, and names in brackets are Lean declarations.
\begin{itemize}
\item T1 [K, \texttt{T1}]: the blocks of terms with 0, 1 and 2 items are disjoint and avoid 36.
\item T2 [K, \texttt{T2}, \texttt{T2\_dead}]: 13 is the only two-item block that revives a dead
application.
\item T3 [K, \texttt{T3a}--\texttt{T3c}]: the table entries equal to 13, 2 and 1 are determined.
\item T4 [K, \texttt{T4}]: 36 absorbs on the left.
\item T5 [K, \texttt{T5}, \texttt{T5\_dies}]: a block recurrence of case (b) of Thm.~VIS has solutions
only at $j=1$.
\item T7 [K, \texttt{T7}, \texttt{T7\_hnfless}]: $U_e\,(S\,U_e)$ has no hnf (folds 33, 31).
\item T8 [K, \texttt{T8}]: folds and absence of hnf for the counters on $X$.
\item T9 [K, \texttt{T9}]: $X$ applied to a right comb has an hnf.
\item T10 [K, \texttt{SELF\_eq}, \texttt{self\_dead}]: the self-dead blocks $\Self$.
\item T11 [K, \texttt{DEAD2\_eq}, \texttt{two\_dead}]: the 2-dead blocks.
\item T12 [K, \texttt{STK\_eq}, \texttt{T12}]: every stack block is self-dead.
\item T13 [K, \texttt{T13}]: the one-item 2-dead block is 27, and $S\,W\,t$ over it has fold 31.
\item (No T14. \S16.2 has an unnamed information line of \texttt{s16tab.log}: $\beta(0,\sigma)$ for
self-dead $\sigma$.)
\item T15 [K for membership, \texttt{T15}, \texttt{tee\_self\_dead}; the values 34, 31 are [C]]: if $a$
is the stack of a normal pair, $S\,a\,(S\,a)$ is self-dead.
\item T16 [K, \texttt{TH1}, \texttt{TH2}, \texttt{T16}, \texttt{T16\_dead}]: the folds of $\Hc$ lie in
$\Theta$, and two-item members are dead against all members.
\item T17 [K, \texttt{CB}, \texttt{IB}, \texttt{T17}, \texttt{T17\_dead}; the equality
$IB=\{12,22,30\}$ is [C]]: comb and inert folds; two-item members are dead against inert terms.
\item T18 [C]: $\fold(Y_m)=31$, and row 31 is 36 on the comb folds (Prop.~NPF).
\item T19 [C]: if $f\,x$ has no hnf and $f$, $x$ have hnfs, then $f$ is self-dead (domain as repaired,
W22.3).
\item T20 [C]: applications over $\Theta$ have fold 31 or 34; rows 31 and 34.
\item T21 [C]: $X\,(S\,W)$, $U\,D(W)$ and (as repaired, W24.2) $Y\,W$ have no hnf for $W=S\,c\,Y$,
$Y\in\{X,U\}$.
\item T22 [C]: the AD folds are exactly $\{31,34\}$ (as repaired, W25.9), closed under $S\,p\,q$.
\item T23 [C, information; \texttt{s26tab.py}, not used in proofs]: the folds of the landing heaps
over a self-dead two-item $Z$. $S\,Z\,v$ is AD; the fold of $S\,c\,Z$ depends only on $\fold c$, is not
self-dead iff $\fold c\in CB$ (then in $IB$), is 25 (self-dead, not AD) iff $\fold c\in\{3,6,7,12\}$,
and is AD otherwise; $\beta(\beta(0,p),q)\in\Self$ for $p\in\Self$.
\item T24 [C, \texttt{s26tab.py}; to port, W26.9]: for $\zeta\in\Self$, $\beta(\zeta,y)=36$ for every
$y\in\mathrm{Ctx}(\zeta)$ (\see{g:wl-ctx}).
\item T25 [C, \texttt{s27tab.py}; to port]: $\beta(f,y)=36$ for every $f\in\Self\cup\{36\}$ and every
$y\in\mathrm{Ctx}_\cup$, with T24 as a special case (\see{g:wl-cdplus}).
\end{itemize}
\end{entry}

\begin{entry}{g:wl-tape}{tape class $\mathcal T$, Conj.~LOAD}{2026-11-26}
\src{LOCK1-INV \S I1 (as repaired, WI.5); the definition is the code's,
checks/s4/lock1-inv/inv\_data.py (t12) and invlib.py (pop, lam)}
\status{open [Conj]; computed [C]: 0 failures in 107{,}264 refill landings across 21{,}275 runs
carried through 30 or 40 moves}
A class of normal pairs: $(A,W)\in\mathcal T$ if, for some right suffix $\sigma$ of $W$, the stack $A$
lies on the pop chain $\lambda(\sigma)=(\sigma,\pop\sigma,\pop^2\sigma,\dots)$ and every term on
$\lambda(\sigma)$ has its fold in $\Self$; the chain stops before $T=S\,S\,S$ or where $\pop$ is
undefined (\see{g:pop}; not the left spine of \see{g:left-spine}). Conjecture LOAD: in a run whose
teeth all normalize, every refill after the first lands in $\mathcal T$. A single move with a
normalizing tooth can leave $\mathcal T$ (487 exits, all in runs that die), so $\mathcal T$ is not an
inductive invariant, and a proof of LOAD needs more than the current pair and its tooth.
\end{entry}

\begin{entry}{g:wl-lock1-ct-ntd}{Theorem~4 of LOCK1-structure (Lock~1 $\Leftrightarrow$ CT $\wedge$ NTD)}{2026-11-26}
\src{LOCK1-structure \S3}
\status{proved [P], classical}
Theorem. Lock~1 holds iff CT and NTD both hold; if Lock~1 fails, there is a spine pair or a bad pair
whose descent turns infinitely often. It splits Lock~1 into two of its consequences rather than
reducing it to a known statement.
\end{entry}

\begin{entry}{g:wl-theorem-c}{Theorem C, counter frame (PF, PFc)}{2026-11-26}
\src{LOCK1-CTB \S9.2, \S10.5}
\status{proved [P]; not in the final chain}
A PF frame is $y\hred S\,(S\,S)\,\delta$ (Lemma~F), and a counter frame PFc adds $\delta=Q^jG$.
Theorem~C: one PFc move after the first B-move of an S-run gives $a'_n\conv a_n$, so $q\conv p_K$.
\end{entry}

\begin{entry}{g:wl-theorem-c-cong}{Theorem C of LOCK1-CONG ($\approx$ is a congruence)}{2026-11-27}
\src{LOCK1-CONG \S3 (as repaired, WC.1--WC.8)}
\status{proved [K]; read as equality of $t^\infty$, given S3's identification of $\mathrm{NF}$ prefixes with the limit}
\lean{\texttt{Congruence.eqv\_app}, \texttt{Congruence.eqv\_ctx}}
Theorem (not Theorem~C of LOCK1-CTB, \see{g:wl-theorem-c}). If $t\approx t'$ and $u\approx u'$ for closed
$S$-terms, then $t\,u\approx t'\,u'$, and $\approx$ is preserved by every context. Hence terms with equal trees in
every context are exactly the $\approx$-pairs, and Lock~1 $\Leftrightarrow$ CL1 (\see{g:wl-cl1}); in Lean,
$\approx$ is \texttt{Eqv}, equal $\mathrm{NF}$ prefixes at every depth, on pure $S$-terms.
\end{entry}

\begin{entry}{g:wl-theorem-k}{Theorem K, Theorem K$_w$}{2026-11-26}
\src{LOCK1-CTB \S3, \S8.0}
\status{proved [P]; on $p_K$'s run without the block model (WB6.5), off the run given (S1) and an exact
check}
Theorem. The run of $p_K$ is explicit: every move is a one-round B-move whose frame term is a two-item
hnf, no tooth has an hnf, and no move is exposed; so exposure need not recur in a run with infinitely
many non-normalizing teeth. K$_w$ is the same run with any closed heap $w$ in place of the last $S$.
\end{entry}

\begin{entry}{g:wl-theorem-kh}{Theorem KH}{2026-11-26}
\src{LOCK1-CTB \S7.2}
\status{proved [P], as repaired (WB2.1)}
\appears{S11 as Lemma~11}
Theorem. If $q$ has $p_K$'s teeth, then from its first B-move on every move is a B-move with heap $S$
or with heap convertible with $\tau_1(Y_n)$, with $k=0$ except possibly once. So $q$ is an S-run or
falls in case (H1) (\see{g:wl-h1}).
\end{entry}

\begin{entry}{g:wl-theorem-l}{Theorem L}{2026-11-26}
\src{LOCK1-CTB \S8.2, \S12.4}
\status{proved [P]; without the block model after WB6.5}
\appears{S11 as Lemma~13}
Theorem. In an S-run with $p_K$'s teeth, $b'_n\conv b_n$ for every $n$ past the first B-move: iterated
S-cancellation peels the counter one $Q$ at a time. This reduces (ii) to (ii$'$).
\end{entry}

\begin{entry}{g:wl-theorem-re}{Theorem RE, $\mathrm{NTD}_Z$}{2026-11-26}
\src{LOCK1-NTD-probe \S5.2}
\status{proved [P], as repaired (WN5.1)}
Theorem. If $t\approx u$ and an invariant $\mu$ of the right end differs on $t$ and $u$, the pair is bad
and NTD fails along the rightmost branch of $t^\infty$ or, after one count, of a child's subtree. Hence
$\mathrm{NTD}_Z$ implies that $\mu$ is a function of the tree.
\end{entry}

\begin{entry}{g:wl-theorem-v}{Theorem V (visible holes are injective)}{2026-11-26}
\src{LOCK1-INJ \S1}
\status{proved [P] (a corollary of Fact~C, WI1.3)}
Theorem. A context $C[v]$ in which a copy of $v$ is reachable as $\mathrm{comp}_\alpha$ without any run
reading $v$ is injective: $C[z]\conv C[z']$ implies $z\conv z'$. Read positions can fail: with
$T=S\,S\,S$ and $a=S\,T\,(S\,T)$, $a\,T\conv a\,(S\,T\,T)$ although $T\not\conv S\,T\,T$ [C].
\end{entry}

\begin{entry}{g:wl-thin-perfect}{thin and perfect pairs}{2026-11-26}
\src{LOCK1-structure \S2 (Theorems~3, 3$'$)}
\status{proved [P], classical}
Theorem. A bad pair either has a thin node, below which the bad set is a single infinite path with all
off-path child pairs convertible, or is perfect: every bad node branches. A thin pair has one context
$E_n$ of each depth with $p\conv E_n[p_n]$ and $q\conv E_n[q_n]$ (Theorem~3$'$).
\end{entry}

\begin{entry}{g:wl-tz}{TZ}{2026-11-26}
\src{LOCK1-CT \S23.3, \S24.1, \S24.2, \S25.1, \S26.3 (Cor.~TZK1), \S26.4, \S27.3 (Cor.~NTZ1)}
\status{refuted [P] (Cor.~NTZ1); given (S1), T8, T12 [K] (W27.1); true for every $b\ne S$ and for
$b=S$ over $Z\in\mathcal S$ (Cor.~TZK1, Prop.~TZX$_{\mathcal S}$)}
Conjecture, now false (T-round landings). If $(S\,T\,b,Z)$ is a normal pair whose first move is a
refill with $k=1$ and all of whose teeth normalize, then $(Z,\nf(b\,Z))\in\Gc^{x}$; the T-start $(X,Z)$
of Prop.~NTZ has all teeth normalizing and lands outside $\Gc^{x}$. Its timed form $\mathrm{TZ}_e$
replaces it (\see{g:wl-tze}, \see{g:wl-ntz}).
\end{entry}

\begin{entry}{g:wl-tze}{$\mathrm{TZ}_e$, $\mathrm{TZX}_e$, Prop.~TZE}{2026-11-26}
\src{LOCK1-CT \S27.4}
\status{open [Conj]; no candidate against $\mathrm{TZX}_e$ among 1,500,410 distinct open-class starts
[C]; Prop.~TZE proved [P], given (S1), T10, T12, T16, T17 [K], T20, T21, T24 [C]}
Conjecture $\mathrm{TZ}_e$: if $(S\,T\,b,Z)$ is a normal pair whose first move is a refill with $k=1$
and all of its teeth normalize, the run of the landing $(Z,\nf(b\,Z))$ has a state in $\Gc^{x}$ by the
state right after its own first refill. $\mathrm{TZX}_e$ is the case $b=S$, $Z$ two-item, self-dead,
$Z\notin\Hc^{x}\cup\mathcal S$, and Prop.~TZE says $\mathrm{TZ}_e\Leftrightarrow\mathrm{TZX}_e$; Prop.~NTZ
does not contradict it (entry at index 2, first refill at index 3).
\end{entry}

\begin{entry}{g:wl-tzk}{Prop.~TZK, Cor.~TZK1}{2026-11-26}
\src{LOCK1-CT \S26.3}
\status{TZK proved [P], given (S1), T12 [K], T24 [C]; TZK1 proved [P], given (S1), T10, T12, T16,
T17 [K], T20, T21, T24 [C] (Wizard: sound; bounded death move with its fold-36 witness as repaired,
W26.7)}
Proposition TZK. For two-item self-dead $Z\notin\Hc^{x}$ and $b\in\Kc\setminus\{S\}$, some tooth of
$(Z,\nf(b\,Z))$ has no normal form, after at most two T-refills (two only if $b=T$); this closes (N1)
for $b\ne S$ and the (N2) recursion. Corollary TZK1: TZ holds for every $b\ne S$, and for $b=S$ it is
TZN over the landing $(Z,S\,Z)$, which Prop.~NTZ refutes (\see{g:wl-ntz}).
\end{entry}

\begin{entry}{g:wl-tzn}{TZN}{2026-11-26}
\src{LOCK1-CT \S25.1 (Cor.~TS1), \S25.4, \S25.5, \S26.3 (Prop.~TZK), \S26.4}
\status{proved for $b\ne S$ [P], given (S1), T12 [K], T24 [C] (Prop.~TZK); for $b=S$ proved over
$Z\in\mathcal S$ (Prop.~TZX$_{\mathcal S}$, same inputs) and refuted [P] over $Z\notin\mathcal S$
(Prop.~NTZ; given (S1), T8, T12 [K], W27.1)}
$\mathrm{TZN}(Z)$, for two-item $Z\notin\Hc^{x}$: for every $b\in\Kc$ some tooth of $(Z,\nf(b\,Z))$
has no normal form. For $b\in\Kc\setminus\{S\}$ the landing run makes only pops and at most two
T-refills before a tooth dies, and Prop.~TZA holds vacuously (W26.5). Kill test [C] (\S26.6): 5,806,281
distinct starts on engine~1 (773,885 also on engine~2), all certified dead, with no budget stop and no
complete run.
\end{entry}

\begin{entry}{g:wl-tzx}{Prop.~TZX$_{\mathcal S}$, Conj.~TZX$'$}{2026-11-26}
\src{LOCK1-CT \S26.4}
\status{TZX$_{\mathcal S}$ proved [P], given (S1), T12 [K], T24 [C] (Wizard: sound); TZX$'$ refuted
[P] (Prop.~NTZ), given (S1), T8, T12 [K] (W27.1)}
Proposition TZX$_{\mathcal S}$. For self-dead $Z\in\mathcal S\setminus\Hc^{x}$ of depth $r$, some tooth of
$(Z,S\,Z)$ has no normal form, within $2r+2$ teeth. Conjecture TZX$'$, now false: the same for two-item
self-dead $Z\notin\Hc^{x}\cup\mathcal S$; of 1,500,410 such starts all die certified except the start
$Z$ of Prop.~NTZ, whose teeth all normalize (\see{g:wl-ntz}, \see{g:wl-tze}).
\end{entry}

\begin{entry}{g:wl-frakB}{the typing $\Bk$, passive slots, Lemmas CL and B}{2026-11-26}
\src{LOCK1-CTB \S11.1--11.2, \S12.4}
\status{proved [P]}
\appears{S11, \S5 ($\Bk$); Lemma~CL as Lemma~16, Lemma~B as Lemma~17}
\lean{\texttt{S11.InB}}
The notes' typing of the world of $p_K$, which S11 publishes as $\Bk$ (\see{g:frakB}); trailing
arguments are passive slots, never moved into function position. Lemma~CL: a set closed under parts
and under the contracta of its three-argument head-run terms has its reducts closed under reduction
and subterms. Lemma~B: $\Bk$ is such a set.
\end{entry}

\begin{entry}{g:wl-vb}{VB (Prop.~VBP)}{2026-11-26}
\src{LOCK1-CT \S23.3 (Conj.), \S24.3, \S26.5 (Prop.~VBP)}
\status{proved [P], given (S1), T12 [K], T24 [C], and T1--T4 [K] through Lemma~XB (droppable, W18.6)
(Prop.~VBP; Wizard: sound)}
Proposition VBP (formerly Conj.~VB). If $V$ is normal with two items, $V\,D(S\,V)$ has no hnf and all
teeth of $(X,S\,V)$ normalize, then $V\in\Hc^{x}_2$: the heap $D(S\,V)\in\mathcal D(V)$ is
context-dead, so Cor.~CDS puts $V\in\mathcal S$ with bottom $X$ or $U$. The reverse statement is not
``the converse'' (W26.1): $V=S\,(S\,X)\,U\in\Hc^{x}_2$ dies at move~2 of $(X,S\,V)$; the kill test of 199,472
distinct starts found all dead in the predicted shape [C].
\end{entry}

\begin{entry}{g:wl-w}{(W), Theorem W, Theorem W6}{2026-11-26}
\src{LOCK1-CTB \S10.3, \S11.3}
\status{(W) proved [P] (Theorem W6); Theorem W as repaired (WB5.1, WB5.2)}
(W): every solution from $k$ that is a subterm of a reduct of $p_K$ satisfies $s\,S\conv a_k\,S$.
Theorem~W: (ii$'$) $\Leftrightarrow$ (W); Theorem~W6 proves (W): every solution in $\mathrm{Red}(\Bk)$ is
convertible with $a_k$. These are not the review labels W$n$.$k$.
\end{entry}

\begin{entry}{g:wl-xu-plus}{$\mathrm{XU}^{+}$}{2026-11-26}
\src{LOCK1-CT \S18.3}
\status{open [Conj]; $\mathrm{XU}^{+}\Rightarrow\mathrm{REC}_T$ proved [P] (Cor.~XU)}
Conjecture. In a normalized run all of whose teeth normalize, $X$ or $U$ over a self-dead heap recurs:
for every $n_0$ there is $n>n_0$ with $A_{n-1}\in\{X,U\}$ and $W_{n-1}$ self-dead.
\end{entry}

\bibliographystyle{plainnat}
\bibliography{refs}
\end{document}
